% FORGE, ICLR 2027 conference manuscript.
% Quantitative inputs are byte-identical snapshots of the recorded source artifacts.
% See iclr_provenance.json before refreshing generated tables, figures, or bibliography files.

\documentclass{article}
\usepackage{iclr2027_conference,times}
% Anonymous review mode; enable \iclrfinalcopy only for an accepted camera-ready version.

\usepackage[utf8]{inputenc}
\usepackage[T1]{fontenc}
\usepackage{nicefrac}
\usepackage{microtype}
\usepackage{amsmath,amssymb,amsthm,mathrsfs}
\usepackage{algorithm}
\usepackage{algpseudocode}
\usepackage{graphicx}
\usepackage{pgfplots}
\usepackage{pgfplotstable}
\pgfplotsset{compat=1.18}
\usepackage{tikz}
\usetikzlibrary{arrows.meta,calc,positioning,patterns,shapes.arrows,fit,backgrounds}
\usepackage{array}
\usepackage{booktabs}
\usepackage{colortbl}
\usepackage{longtable}
\usepackage[most]{tcolorbox}
\usepackage{float}
\usepackage{placeins}
\usepackage{chemfig}
\usepackage{xcolor}
\usepackage{url}
\usepackage[hidelinks]{hyperref}

% Figures and generated rows are local, hash-recorded copies of the source manuscript inputs.
\graphicspath{{./}}

\newtheorem{theorem}{Theorem}[section]
\newtheorem{proposition}[theorem]{Proposition}
\newtheorem{lemma}[theorem]{Lemma}
\newtheorem{corollary}[theorem]{Corollary}
\theoremstyle{definition}
\newtheorem{definition}[theorem]{Definition}
\newtheorem{assumption}[theorem]{Assumption}
\newtheorem{invariant}[theorem]{Implementation invariant}

% Formal statements use a quiet lavender field to separate assumptions and guarantees from the
% surrounding argument. The box may break across pages; proofs intentionally remain unboxed.
\definecolor{forgeformalbg}{HTML}{F0EFFF}
\definecolor{forgerow}{HTML}{E9E7FB}
\tcbset{
  forge formal block/.style={
    enhanced jigsaw,
    breakable,
    colback=forgeformalbg,
    colframe=forgeformalbg,
    boxrule=0pt,
    arc=0pt,
    outer arc=0pt,
    boxsep=0pt,
    left=8pt,
    right=8pt,
    top=7pt,
    bottom=7pt,
    before skip=8pt,
    after skip=8pt
  }
}
\tcolorboxenvironment{definition}{forge formal block}
\tcolorboxenvironment{assumption}{forge formal block}
\tcolorboxenvironment{proposition}{forge formal block}
\tcolorboxenvironment{theorem}{forge formal block}
\tcolorboxenvironment{lemma}{forge formal block}
\tcolorboxenvironment{corollary}{forge formal block}
\tcolorboxenvironment{invariant}{forge formal block}

% Unresolved experimental results remain visually explicit until a verified artifact supplies
% the value. Never replace one of these placeholders from an unverified log or planning document.
\newcommand{\RESULTPENDING}[1]{\textbf{\textcolor{red}{[PENDING: #1]}}}
% Generated from the hash-pinned completed-evidence contract. Do not edit numerical values here.
% Missing measurements remain explicit in a final-paper narrative.
\definecolor{forgepending}{HTML}{9B4A12}
\newcommand{\PendingCell}{\textcolor{forgepending}{\textsf{---}}}
\newcommand{\PendingValue}[1]{\textcolor{forgepending}{\textsf{[#1]}}}
\newcommand{\ResultPanel}[2]{\fcolorbox{black!25}{black!2}{\parbox[c][2.1cm][c]{0.43\linewidth}{\small\textbf{#1}\par\medskip #2\par\smallskip\PendingValue{result}}}}


\input{generated/completed_evidence_macros.tex}
% Superseding Table 1 values from the matched final conditioned, null and cyclic evaluations.
\input{generated/iclr_table1_macros.tex}
% Superseding Table 3 values from the final graph Transformer population and frozen route index.
\input{generated/iclr_table3_route_macros.tex}
% Superseding Table 10 values from the final method-blind common-Ugi route-evidence adjudication.
\input{generated/iclr_table10_macros.tex}
% Every computed value cited in prose is regenerated from the same final evidence as its table.
\input{generated/iclr_prose_results_macros.tex}

% Appendix generated-structure atlas. The hash-pinned row file carries structures and SMILES only
% and defers every length, rule and type style to the definitions below, so the appendix can be
% retuned without reissuing a provenance artifact.
\definecolor{atlasrule}{gray}{0.74}
\definecolor{atlasink}{gray}{0.42}
\newlength{\atlasgraphwidth}
\newlength{\atlasprecursorwidth}
\newlength{\atlasgraphheight}
\newlength{\atlasprecursorheight}
\newlength{\atlasrowheight}
\newlength{\atlasrankwidth}
\newlength{\atlassmilesskip}
% Every column is pinned to a stated width, so the alignment cannot be widened past \textwidth by a
% heading or an index label. These must sum with the three gutters and the divider in the column
% specification: 5mm + 1.8mm + 0.4549 + 2.8mm + 0.5pt + 2.8mm + 0.4549 = \textwidth.
\setlength{\atlasrankwidth}{5mm}
\setlength{\atlasgraphwidth}{0.4549\textwidth}
\setlength{\atlasprecursorwidth}{0.4549\textwidth}
\setlength{\atlasgraphheight}{2.45cm}
\setlength{\atlasprecursorheight}{3.80cm}
% Rows are struck to a common height that exceeds both panels: the surplus is the air between
% structures, and holding it inside the row keeps the column divider unbroken and keeps the
% multipage break in the same place for every row rather than following a wrapped SMILES string.
\setlength{\atlasrowheight}{4.00cm}
\setlength{\atlassmilesskip}{2.4mm}
\newcommand{\atlassmilesstyle}{\scriptsize\ttfamily\color{atlasink}}
\newcommand{\atlasrank}[1]{%
  \makebox[\atlasrankwidth]{\raisebox{-0.5\height}{\footnotesize\color{atlasink}#1}}}
\newcommand{\atlashead}[2]{\makebox[#1]{\footnotesize\scshape #2}}

\title{FORGE: Reaction-Guided Generative Design of Ionizable Lipids}

\author{Anonymous authors\\Paper under double-blind review}

\begin{document}

\maketitle
\raggedbottom

\begin{abstract}
Ionizable lipids are central components of lipid nanoparticles for RNA delivery. Their discovery
often relies on synthesizing and screening libraries assembled from predefined reactants, which
restricts the precursor structures available for design. Whole-molecule generation can extend
this search, but verifying assembly requires precursor assignments and reaction information.
We introduce FORGE (\textbf{F}low-matched, \textbf{O}pen-ended,
\textbf{R}oute-resolved \textbf{G}eneration and \textbf{E}xploration), which uses known reaction chemistry to design ionizable lipids without restricting
generation to a fixed reactant catalogue. Each reaction program specifies precursor roles,
reaction-core positions, transformation order and coarse molecular architecture. FORGE generates
the complete molecule at atom-and-bond resolution without component identifiers, stored precursor
graphs or fragment tokens. Reverse decomposition and exact forward replay verify the implied
assembly. We evaluate one shared model across 22 assembly families, including multicomponent
condensations, repeated additions, ring openings and acylation reactions. Across
\PendingValue{independent training seeds}, verified exact level-one (exact-L1) yield is
\PendingValue{macro mean and seed SD}\%, compared with \PendingValue{matched null yield}\%
for whole-molecule generation followed by post-hoc verification. Structural checks, reference-lipid
distributions and component-level diversity distinguish reaction consistency from molecular quality.
The model generates \PendingValue{distinct component-novel yield}\% exact-L1 products with
recovered precursor identities absent from the training catalogues. Recursive synthesis assessment
resolves complete precursor-preparation and terminal-material dossiers for
\PendingValue{complete dossier yield}\% of requests. These results
\PendingValue{conclusion supported by conditioning, quality and route comparisons}.
\end{abstract}

\section{Introduction}

Lipid nanoparticles (LNPs) deliver siRNA, mRNA vaccines and genome-editing reagents
\citep{adams2018patisiran,polack2020safety,baden2021efficacy,gillmore2021crispr}, with ionizable
lipids supporting nucleic-acid complexation, particle assembly and endosomal escape
\citep{cullis2017lipid,hou2021lipid}. Discovery of these lipids commonly relies on combinatorial
synthesis and screening, in which modular precursors are assembled into libraries and the resulting
formulations are evaluated \citep{akinc2008combinatorial,love2010lipidlike,li2024accelerating,xu2024agile}.
This approach uses established reactions to combine known building blocks, which limits the lipids it can generate.

Generative methods incorporate assembly chemistry at different stages of molecular design.
SynNet, RGFN and SynFlowNet retain reaction routes by selecting reactants from a supplied inventory
and applying predefined transformations \citep{gao2022amortized,koziarski2024rgfn,cretu2025synflownet},
a strategy also used by lipid-specific synthesis-DAG and tree-search methods
\citep{ou2024deep,zhao2024generative}. SynCoGen couples precursor design to assembly by jointly
generating building-block graphs, reaction assignments and atomic coordinates
\citep{rekesh2025syncogen}. Whole-molecule models such as DeFoG and GenMol generate structures
without a prescribed precursor inventory \citep{qin2025defog,lee2025genmol}, leaving assembly and
synthesizability to separate assessment \citep{gao2020synthesizability}.

FORGE learns a conditional distribution over complete lipid graphs from reaction programs that
specify roles, reaction-core positions, transformation order and coarse architecture. It receives
no precursor identifiers, stored component graphs or fragment tokens. Reverse decomposition and
exact forward replay verify each generated assembly. A bounded recursive search then assesses
whether the implied precursors have preparation routes supported by reaction and supplier evidence,
returning a computational dossier or an explicit unresolved outcome.

Our contributions are:
\begin{enumerate}
    \item \textbf{Reaction-program-conditioned whole-molecule generation.}
    We developed a conditional discrete flow that generates complete lipid graphs with precursor
    identities left open, coupled to reverse decomposition and exact forward replay for assembly
    verification.
    \item \textbf{Matched evaluation across 22 assembly families.}
    We compare one shared model with matched whole-molecule generation followed by post-hoc
    verification, conditioning controls and component-based generators. Family-wise assembly,
    decomposition precision and coverage, and held-component recovery quantify the contribution
    of reaction-program conditioning beyond the training inventory.
    \item \textbf{Structural quality and synthesis evidence.}
    We assess headgroup retention, ring structure, chemical alerts, lipid-reference realism and
    product and precursor diversity alongside exact assembly. Recursive route assessment records
    precursor-preparation evidence, terminal-material support and unresolved outcomes for the
    generated identities, separating L1 consistency from complete synthesis dossiers.
\end{enumerate}

\section{Preliminaries and problem formulation}
\label{sec:preliminaries}

\subsection{Ionizable lipids and molecular architecture}
\label{sec:lipid-background}
Ionizable lipids combine protonatable amine-containing headgroups, linkers and hydrophobic tails
\citep{semple2010rational,jayaraman2012maximizing}. Protonation supports nucleic-acid association
during acidic formulation and membrane interactions involved in endosomal release
\citep{xu2024agile}. LNPs commonly combine an ionizable lipid with helper lipids, cholesterol and
PEGylated lipids. Lipid structure, stereochemistry and formulation composition affect delivery
\citep{han2021ionizable,hatit2023nanoparticle,cheng2020selective}; molecular constitution and
formulated-particle performance are therefore distinct evaluation levels.

\subsection{Reaction families and modular assembly}
\label{sec:reaction-background}
We study 22 registered assembly families (Table~\ref{tab:reaction-family-background}).
These include multicomponent condensations, repeated amine additions, ring-opening reactions,
esterification and amidation. Ugi 3-CR combines an amine, aldehyde and
isocyanide into an $\alpha$-amino amide \citep{xu2024agile}. Repeated aza-Michael addition attaches
acrylate-bearing components to amine sites \citep{li2023combinatorial}; repeated reductive amination
couples aldehyde-bearing components to amines through condensation and reduction
\citep{xue2024barcoding}. The Ugi variant has no carboxylic-acid reactant component and uses an
acidic phosphorus catalyst; its ester functionality comes from the aldehyde precursor.
Repeated assemblies require both transformation order and participating sites. Precursor-origin
boundaries can span functional headgroup, linker and tail regions. These distinctions motivate
role-assigned reaction programs (Appendix~\ref{app:chemistry}, Table~\ref{tab:reaction-family-background}).

\subsection{Molecular graphs and reaction programs}
\label{sec:mathematical-objects}
A lipid constitution is a nonempty, connected, simple molecular graph
$G=(V,E,\ell_V,\ell_E)$, with atom labels $\ell_V:V\to\Sigma_{\rm atom}$ and bond labels
$\ell_E:E\to\Sigma_{\rm bond}$ in finite vocabularies. Bond order is an edge label.
We compare graphs up to label-preserving isomorphism after removing atom-map and stereochemical
labels, retaining formal charges and bond orders. We write $G$ for a graph or its constitutional
isomorphism class when the distinction is immaterial. The bounded graph space
$\mathfrak G_{N,K_{\rm cl}}$ contains these classes with $1\leq|V|\leq N$ and cycle rank
$\beta(G)=|E|-|V|+1\leq K_{\rm cl}$. This is a structural support definition; molecular validity
additionally requires valence and sanitization checks. Here $N=N_{\max}=254$ and $K_{\rm cl}=12$.

A reaction program specifies the assembly constraints,
\begin{equation}
    P=(c,d,\mathbf r,\mathbf k,M).
    \label{eq:reaction-program}
\end{equation}
Here $c$ identifies a registered assembly chemistry with its admitted transformation sequences,
$d$ is the requested depth, $\mathbf r$ assigns precursor roles to
atom positions, and $\mathbf k$ marks reaction-core positions. The four morphology fields in $M$
record each role's exterior atom count, junction budget, closure count and core-attachment count.
These quantities constrain coarse architecture without specifying precursor identities. Precursor
origin and core membership are separate: an atom can belong to a precursor and participate in the
core; a template-introduced atom belongs to the core without belonging to a precursor.
Appendix~\ref{app:program-encoding} defines the morphology counts.

\subsection{Exact assembly, novelty and route evidence}
For registered chemistry $c$, a precursor input $\mathbf C\in\mathfrak C_c$ is an ordered tuple
of constitutions, roles and transformation steps. Its registered forward products within the
bounded support form $\mathcal F_c(\mathbf C)\subseteq\mathfrak G_{N,K_{\rm cl}}$.
Reverse search returns a finite set $\mathcal D_c(G)$ of precursor inputs and assembly traces
$(\mathbf C,\tau)$. Exact forward replay retains
\begin{equation}
    \widehat{\mathcal E}_c(G)=
    \left\{(\mathbf C,\tau)\in\mathcal D_c(G):G\in\mathcal F_c(\mathbf C)\right\}.
    \label{eq:main-exact-traces}
\end{equation}
Membership uses constitutional identity. A valid graph passes verified exact level-one assembly
(L1) when $\widehat{\mathcal E}_c(G)\neq\varnothing$. Bounded reverse search can miss assemblies,
and recovered traces need not match every sampled depth, role-position or morphology field in $P$.
Assembly existence and conditioning adherence are distinct checks, as are decomposition coverage
and replay precision (Appendix~\ref{app:verification}).

A recovered precursor is component-novel when its constitution is absent from its role's catalogue.
Strong product-level catalogue escape requires \emph{every} exact precursor input to contain an
outside component. We measure recovered component novelty, which does not establish that universal
claim (Appendix~\ref{app:catalogue-reachability}). A complete computational synthesis dossier also
requires evidence-qualified precursor preparation (L2) and terminal-material availability (L3).
Incomplete assessments remain unresolved or search-censored; these checks do not establish
experimental synthesis success or delivery activity.

\section{FORGE}
\label{sec:method}

\subsection{Overview}
\label{sec:forge-overview}

FORGE learns the distribution $p_\theta(G\mid P)$ of complete lipid molecules $G$ given a
reaction program $P$. The model receives no precursor identifiers, stored component graphs or
fragment tokens. The program specifies how the molecule is assembled without selecting the
precursor structures. FORGE has four main steps (Figure~\ref{fig:forge-overview}):
\begin{enumerate}
    \item \textbf{Specify the reaction program.} Choose assembly chemistry, roles, reaction-core
    positions, transformation order and morphology, leaving precursor identities open.
    \item \textbf{Generate the whole graph.} A conditional discrete flow generates atom, bond and
    topology states while preserving program-fixed coordinates.
    \item \textbf{Verify exact assembly.} Recover role-assigned precursors and require their forward
    transform to reconstruct the generated molecule.
    \item \textbf{Construct precursor routes.} Search for preparation routes to evidence-qualified
    starting materials and abstain when a complete dossier cannot be constructed.
\end{enumerate}

\begin{figure*}[htbp]
\centering
\resizebox{\textwidth}{!}{\input{figures/forge_algorithm_schematic/schematic_body.tex}}
\caption{\textbf{FORGE generation and route assessment.}
A reaction program fixes assembly roles and core coordinates without selecting precursor identities.
FORGE generates the whole graph, verifies exact assembly replay, and recursively traces its implied
precursors to evidence-qualified materials or an explicit abstention. Colours mark precursor origin;
grey denotes the Ugi template-introduced oxygen. Replay and route closure are computational checks,
not experimental synthesis evidence.}
\label{fig:forge-overview}
\end{figure*}

\subsection{Conditioning on reaction programs}
\label{sec:program-conditioning}
An admissible complete program is sampled from the training-fold prior,
\begin{equation}
    p_\theta(G,P\mid c)=p_{\mathrm{prog}}(P\mid c)\,p_\theta(G\mid P).
\end{equation}
Learned embeddings combine chemistry, depth, precursor role, core status and role-local morphology
at each graph position. The adapter determines active positions and fixed or variable coordinate
masks. Appendix~\ref{app:program-encoding} specifies the embeddings and morphology counts.

\subsection{Graph representation and discrete flow}
\label{sec:graph-flow}
FORGE represents each molecule as a spanning tree with additional bonds for ring closures.
A deterministic encoder maps the molecular graph $G$ and its selected atom-mapped
assembly annotation $\alpha$ to
\begin{equation}
    X_\star := \operatorname{Enc}_P(G,\alpha)
    = (\mathbf a,\boldsymbol\pi,\mathbf b,\mathbf e^L,\mathbf e^R,\mathbf z),
\end{equation}
where $\mathbf a$ records atom states, $\boldsymbol\pi$ identifies each atom's parent in the tree,
and $\mathbf b$ records the tree bonds. The arrays $\mathbf e^L$ and $\mathbf e^R$ identify the
endpoints of additional bonds, whose states are stored in $\mathbf z$.

The padded state space $\overline{\mathfrak X}_{N,K_{\rm cl}}$ is finite. Its valid encodings form
$\mathfrak X_{N,K_{\rm cl}}$: each contains a connected tree and at most $K_{\rm cl}$ distinct
closure edges. The decoder $\operatorname{Dec}$ reconstructs the graph from a valid encoding and
returns $\bot_{\rm dec}$ for a malformed state. Every graph within the specified atom, bond, size
and closure limits has an encoding (Theorem~\ref{thm:bounded-representation}).

During sampling, we also restrict the allowed parents and number of children for each atom.
These constraints limit which valid encodings the sampler can use.

For a given $P$, let $\mathcal I(P)$ index the categorical coordinates, each with finite alphabet
$\mathcal X_i(P)$. Pointer alphabets contain only positions permitted by the layout.
The adapter partitions coordinates into fixed positions $\mathcal F(P)$ with values $x_{P,i}$ and
variable positions $\mathcal V(P)$. The set $\Omega_P$ contains raw states within the allowed alphabets that satisfy the fixed
values. It can contain malformed or chemically invalid states. The source law $q_0^P$ fixes
$\mathcal F(P)$ and uses independent, full-support marginals $q_{0,i}^P$ on variable alphabets.
For a clean endpoint $x_1\in\Omega_P$, the conditional corruption law is
\begin{equation}
    q_{t\mid1}^P(x\mid x_1)
    =\mathbf 1[x_{\mathcal F(P)}=x_{P,\mathcal F(P)}]
    \prod_{i\in\mathcal V(P)}
    \left[(1-t)q_{0,i}^P(x_i)+t\mathbf 1\{x_i=x_{1,i}\}\right],
    \quad 0\leq t\leq1.
    \label{eq:conditional-path}
\end{equation}
Each coordinate interpolates between its source and clean value. Their conditional joint law is a
product, so it generally differs from a linear mixture of the joint source and endpoint laws.

The clean-coordinate predictor $d_{\theta,i}(a\mid x,t,P)$ assigns a probability to each
$a\in\mathcal X_i(P)$. A continuous-time Markov chain (CTMC) uses these predictions to average
conditional one-coordinate rates,
\begin{equation}
 R_{t,i}^\theta(u,v\mid x,P)
 =\sum_{a\in\mathcal X_i(P)}d_{\theta,i}(a\mid x,t,P)
 R^*_{t,i}(u,v\mid a,P),\qquad u\neq v,\quad t<1.
 \label{eq:posterior-rate}
\end{equation}
The conditional rate $R^*_{t,i}$ realizes the corresponding factor in
Equation~\ref{eq:conditional-path}; Appendix~\ref{app:flow-proof} defines it and proves the path
identity. The joint generator permits one coordinate to change at an infinitesimal time.
Coordinate posterior marginals therefore suffice to specify its rates, even when endpoint
coordinates are statistically dependent.

\subsection{Program-conditioned Transformer and training}
\label{sec:model-training}

A graph Transformer predicts clean coordinates from $(X_t,t,P)$. Each layer combines graph
self-attention with parent, child and closure relation biases and cross-attention to program,
depth, role and core tokens. Softly program-routed residual adapters add family-specific parameters.
Qualified atom-mapped traces supervise role and core annotations; unresolved source-trace ambiguity
excludes a record from semantic training.

The source-balanced training law $\mu_{\rm bal}$ covers annotated triples $(G,\alpha,P)$ with equal
mass across the 22 qualified families and uniform sampling over unique constitutions within each
family, with replacement. No second loss reweighting changes this measure. The model combines masked coordinate cross-entropy, role-balanced chemistry, role/core supervision,
and repeat and topology consistency terms. We combined chemistry-task gradients with PCGrad
\citep{yu2020gradient}. Appendix~\ref{app:training} specifies the full objective and training-time
distribution.

To isolate the statistical property of coordinate denoising, define the reference population loss
\begin{equation}
 \mathcal L_{\rm ref}(\theta)
 =\mathbb E_{\substack{(G,\alpha,P)\sim\mu_{\rm bal},\ t\sim\mathcal U(0,1)\\
 X_t\sim q_{t\mid1}^P(\cdot\mid X_\star)}}
 \left[\sum_{i\in\mathcal V(P)}w_i(P)
 \bigl[-\log d_{\theta,i}(X_{\star,i}\mid X_t,t,P)\bigr]\right],
 \qquad w_i(P)>0.
 \label{eq:primary-loss}
\end{equation}
The weights depend only on the program; they can normalize coordinate counts without depending
on the unknown clean category. This reference loss supports the population analysis. It is not
the complete production training objective.

The reference objective has an ideal population interpretation. For a program with positive training
mass, let $q_1^P$ be the law of $\operatorname{Enc}_P(G,\alpha)$ under $\mu_{\rm bal}(\cdot\mid P)$,
and let $\mu_G^P$ be the corresponding graph law.
\begin{theorem}[Ideal conditional generation]
\label{thm:ideal-terminal-law}
Fix such a program $P$, with registered chemistry $c$, finite coordinate alphabets and a
full-support source on each variable alphabet. Assume every training endpoint belongs to $\Omega_P$, decodes to its source
graph, and passes molecular validity and verified exact L1. Assume an unrestricted measurable
coordinate predictor attains a population minimum of Equation~\ref{eq:primary-loss}. Starting from $q_0^P$, exactly simulate the local CTMC
of Equation~\ref{eq:posterior-rate} on $0\leq t<1$. Then its marginal law converges to $q_1^P$ as
$t\uparrow1$, and
\begin{equation}
 \operatorname{Dec}_{\#}q_1^P=\mu_G^P,\qquad
 \Pr_{G\sim\mu_G^P}\!\left[\mathrm{Valid}(G)=1,
 \widehat{\mathrm{L1}}(G,c)=1\right]=1.
\end{equation}
Here $\operatorname{Dec}_{\#}$ denotes the distribution obtained by decoding.
\end{theorem}
Appendix~\ref{app:flow-proof} proves this result by recovering the clean-coordinate posteriors and
verifying the CTMC forward equation. The theorem concerns the reference population objective and the
limit of exact continuous-time sampling. It establishes neither generalization beyond $q_1^P$ nor
correctness of finite training, PCGrad, capped Euler updates or constrained terminal readout.

\subsection{Sampling and exact assembly verification}
\label{sec:sampling-verification}
The numerical sampler uses states $Z_k$, distinct from the ideal CTMC $X_t$. Starting from
$Z_0\sim q_0^P$, it applies
\begin{equation}
 Z_{k+1}=\Pi_P\!\left(\Phi_k^\theta(Z_k,\xi_k)\right),
 \qquad k=0,\ldots,L_{\rm samp}-1,
 \label{eq:numerical-sampling}
\end{equation}
where $\Phi_k^\theta$ is a capped independent-coordinate Euler update, $\xi_k$ contains its random
draws and $\Pi_P$ restores program-fixed coordinates. Fixed coordinates have zero rate and remain
constant (Proposition~\ref{prop:fixed-invariance}); chemical consistency requires verification.
A fresh prediction at $t=1$ drives the constrained argmax readout
$Z_{\rm out}=H_P^\theta(Z_{L_{\rm samp}})$, including family-specific core and reactive-site constraints, or a typed failure.
The checked inference arm augments candidate graphs with qualified structural proposals and selects
under explicit product and precursor diversity constraints. We reassess the final graph and retain
raw outputs as a separate arm; Table~\ref{tab:22-decoder} distinguishes these inference contributions
from the learned generator.
We then decode, sanitize and apply Equation~\ref{eq:main-exact-traces} to each valid graph.
Molecular validation, decomposition, exact replay and root ambiguity are recorded separately.
Appendix~\ref{app:sampling} specifies the numerical operations and their distinction from the ideal
CTMC.

\subsection{Recursive route assessment}
\label{sec:route-assessment}
For a molecule with exactly one verified precursor decomposition $(\mathbf C_T,\tau_T)$,
we search for routes to prepare its precursors within fixed depth and query limits:
\begin{equation}
    Y=\mathsf{Assess}_{\mathfrak B}(G,P,\mathbf C_T,\tau_T)
    \in\mathfrak D\cup\mathfrak F_{\mathfrak B}.
\end{equation}
Here $\mathfrak B$ is the frozen route contract, $\mathfrak D$ the dossier space and
$\mathfrak F_{\mathfrak B}$ its finite typed outcome set. Ambiguous roots cause abstention.
A complete dossier requires molecular validity, exact L1, admitted and verified upstream reactions,
and available terminal materials. Other outcomes are unresolved or search-censored. Assessment
operates on the completed graph without modifying generation. Appendix~\ref{app:route-assessment}
and Algorithm~\ref{alg:forge-inference} specify admission and inference; Appendix~\ref{app:theory}
states the formal guarantees and their assumptions.


\section{Results}
\label{sec:results}

\subsection{Evaluation protocol and matched budgets}
The evaluation compares correctly conditioned, shared-null and cyclically mislabelled models under
matched architecture, training exposure, checkpoint selection and inference budgets. The evaluation
population contains \PendingValue{held-out requests per family per seed} requests per family across
\PendingValue{independent training seeds} independent training seeds. Component-disjoint held-out
products are excluded from training, decoder development and checkpoint selection. Verified exact-L1
yield counts attempts passing molecular validation, decomposition and exact forward replay, with
\emph{all} attempts in the denominator. Post-hoc rejection cannot increase the proposal's per-attempt
acceptance probability (Proposition~\ref{prop:posthoc-budget}). We report means and sample standard
deviations across training seeds, distinguishing retraining variation from conditional Monte Carlo
variation. Raw and final yield count one output per request before and after the declared inference
procedure; failed searches and abstentions remain in the denominator. Appendix~\ref{app:22-results}
specifies the splits, source domains, proposal counts and computational budgets.

\subsection{Shared generation across 22 assembly families}
We first ask whether one conditional flow can cover 22 assembly families while maintaining
performance on the matched three-family comparison. The equal-family exact-L1 yield is
\PendingValue{raw macro L1} for raw readout and \PendingValue{final macro L1} after checked inference
(Table~\ref{tab:22-l1}). The lowest family yield is \PendingValue{lowest family L1}, with
\PendingValue{between-family pattern}. Decomposition coverage, candidate precision and ambiguity
are \PendingValue{decomposition outcomes}; replay of accepted traces is
\PendingValue{accepted-trace replay}. These quantities distinguish an inability to recover a valid
assembly from an inconsistent returned trace. Fixed-coordinate violations and support overflows
occur in \PendingValue{coordinate and support failures} of \PendingValue{total requests} requests.
On common requests for Ugi 3-CR, aza-Michael addition and reductive amination, the change relative to
the three-family model is \PendingValue{three-family bridge effect}
(Table~\ref{tab:22-controls}).

\begin{table}[H]
\caption{\textbf{Exact assembly across 22 ionizable-lipid reaction families.}
Values are mean $\pm$ sample standard deviation across independent training seeds, with all
requests in the yield denominator. Raw and final outputs share request identities; final includes
checked completion, repair and selection. Null and cyclic arms use matched training and inference
budgets. Decomposition coverage and candidate precision are distinct from replay of accepted traces.
Per-seed counts, replay invariants and ambiguity: Appendix~\ref{app:22-results}.}
\label{tab:22-l1}
\centering\footnotesize\setlength{\tabcolsep}{3pt}\renewcommand{\arraystretch}{1.12}
\begin{tabular}{@{}>{\raggedright\arraybackslash}p{.30\textwidth}ccccc@{}}
\toprule
Family & \shortstack{Raw\\L1 (\%)} & \shortstack{Final\\L1 (\%)} & \shortstack{Shared\\null (\%)} & \shortstack{Cyclic\\label (\%)} & \shortstack{Decomp. cov. /\\candidate prec.} \\
\midrule
A3 amine--aldehyde--alkyne & \PendingCell & \PendingCell & \PendingCell & \PendingCell & \PendingCell \\
Acid--epoxide diester & \PendingCell & \PendingCell & \PendingCell & \PendingCell & \PendingCell \\
AEMA & \PendingCell & \PendingCell & \PendingCell & \PendingCell & \PendingCell \\
Aldehyde Ugi 3-CR & \PendingCell & \PendingCell & \PendingCell & \PendingCell & \PendingCell \\
Aldehyde Ugi 4-CR & \PendingCell & \PendingCell & \PendingCell & \PendingCell & \PendingCell \\
$\alpha$-Isocyanoester dihydroimidazole & \PendingCell & \PendingCell & \PendingCell & \PendingCell & \PendingCell \\
Amine alkylation & \PendingCell & \PendingCell & \PendingCell & \PendingCell & \PendingCell \\
Amine--epoxide & \PendingCell & \PendingCell & \PendingCell & \PendingCell & \PendingCell \\
Aryl reductive amination & \PendingCell & \PendingCell & \PendingCell & \PendingCell & \PendingCell \\
Aza-Michael acrylamide & \PendingCell & \PendingCell & \PendingCell & \PendingCell & \PendingCell \\
Aza-Michael acrylate & \PendingCell & \PendingCell & \PendingCell & \PendingCell & \PendingCell \\
Disulfide Michael & \PendingCell & \PendingCell & \PendingCell & \PendingCell & \PendingCell \\
Epoxide O-acylation & \PendingCell & \PendingCell & \PendingCell & \PendingCell & \PendingCell \\
iPhos & \PendingCell & \PendingCell & \PendingCell & \PendingCell & \PendingCell \\
Ketone--isocyanide amide & \PendingCell & \PendingCell & \PendingCell & \PendingCell & \PendingCell \\
Ketone Ugi 4-CR & \PendingCell & \PendingCell & \PendingCell & \PendingCell & \PendingCell \\
Maleate & \PendingCell & \PendingCell & \PendingCell & \PendingCell & \PendingCell \\
O-esterification & \PendingCell & \PendingCell & \PendingCell & \PendingCell & \PendingCell \\
Passerini & \PendingCell & \PendingCell & \PendingCell & \PendingCell & \PendingCell \\
Preassembled thiol-yne & \PendingCell & \PendingCell & \PendingCell & \PendingCell & \PendingCell \\
Reductive amination & \PendingCell & \PendingCell & \PendingCell & \PendingCell & \PendingCell \\
Thiolactone & \PendingCell & \PendingCell & \PendingCell & \PendingCell & \PendingCell \\
\midrule
Equal-family macro & \PendingCell & \PendingCell & \PendingCell & \PendingCell & \PendingCell \\
Pooled all requests & \PendingCell & \PendingCell & \PendingCell & \PendingCell & \PendingCell \\
Worst-family result & \PendingCell & \PendingCell & \PendingCell & \PendingCell & \PendingCell \\
\bottomrule
\end{tabular}
\end{table}

\subsection{Contribution of reaction-program conditioning}
To isolate the effect of reaction information during generation, we compare the conditioned
model with the shared-null and cyclic arms under the same attempt budget and post-hoc verifier.
The conditioned-minus-null difference is \PendingValue{conditioning difference} percentage points,
with \PendingValue{conditioning uncertainty and family variation}. The cyclic control retains
precursor roles, reaction-core coordinates and program depth, isolating the contribution of the
reaction-identity label within a structured program. Its difference from correct conditioning is
\PendingValue{identity-label difference}. The information supplied through layouts, masks, cores and
decoder priors is accounted for separately (Table~\ref{tab:22-controls}); the resulting attribution is
\PendingValue{program versus identity-label conclusion}.

\subsection{Decoder and architecture effects across programs}
With model weights fixed, the raw-to-final exact-L1 change is
\PendingValue{decoder L1 difference}, and the change in outputs satisfying the qualified structural
checks is \PendingValue{decoder design difference} (Table~\ref{tab:22-decoder}). The same source
verifier assesses core-constrained readout, component completion, scoped head retention, role-local
atom/bond proposals and final ring-system repair. Original candidates remain available to selection.
Reselection within the original pool changes yield by \PendingValue{selection-only difference};
adding repaired candidates contributes \PendingValue{repair-pool difference}, at a cost of
\PendingValue{additional proposal and source-check cost}. Larger-pool gains and equal-cost contrasts
are reported separately, because additional search and a changed selection objective do not isolate
the effect of a repair rule.

The architecture comparison separates joint denoising from parameter-matched independent precursor
roles (FACT), cross-attention and gradient-conflict control. Exact-L1 and structural-quality
contrasts are \PendingValue{architecture contrasts}, with
\PendingValue{architecture family dependence} (Table~\ref{tab:22-architecture}). These comparisons
separate learned-model effects from changes made at inference; motif agreement with the same TRAIN
components that define a decoder prior is not independent validation.

\subsection{Precursor novelty, diversity and held-component generalization}
We next ask whether FORGE can preserve exact assembly while generating precursor constitutions
absent from training. Verified component-novel yield is \PendingValue{component-novel yield} per
1,000 requests, including \PendingValue{distinct component-novel yield} distinct products
(Table~\ref{tab:22-diversity}). Finite-inventory methods have zero component novelty relative to
their own visible inventory by construction. Component novelty is verifier-relative and does not
establish product-level catalogue escape. Exact product reachability gives
\PendingValue{catalogue reachability outcome}. Role-specific unique counts and Shannon and Simpson
effective counts are \PendingValue{component diversity outcome}, distinguishing a broad precursor
population from repeated use of a few components.

Held-component recovery is \PendingValue{held-component yield} per 1,000 requests, covering
\PendingValue{held-identity coverage} of reserved identities. Source-disjoint and component-disjoint
exact yields are \PendingValue{generalization yields}, with \PendingValue{size and repeat effects}
across molecular size, closure count and repeated-arm strata (Table~\ref{tab:22-generalization}).
Comparisons with external generators use their qualified chemical scope and a common verifier;
their outcome is \PendingValue{external-baseline comparison}. Unsupported native families remain
unassessed. Fresh noise applied to TRAIN-derived requests is a sampling diagnostic, not evidence
of held-out generalization.

\subsection{Structural quality and agreement with lipid references}
Exact assembly and lipid structural quality measure different properties. The fraction satisfying
exact L1 and all applicable design checks is \PendingValue{raw joint design yield} for raw outputs
and \PendingValue{final joint design yield} after inference
(Table~\ref{tab:22-quality-metrics}). Chemical alerts, loss of a required headgroup site and
ring-condition violations account for \PendingValue{structural failure distribution}. These checks
apply to the complete final graph; source-positive controls give
\PendingValue{control disagreement and abstention rates}. Chemistry-context flags remain separate
from confirmed defects. A retained-basic-site query does not establish pKa, and absence of a narrow
alert does not establish physical stability or delivery activity.

\begin{table}[H]
\caption{\textbf{Structural quality and agreement with design conditions.}
Raw and final arms share requests; values are mean $\pm$ sample standard deviation across training
seeds. Controls are source-positive structures within each metric's qualified scope. Applicable and
assessed counts, control disagreements and abstentions accompany conditional rates; unsupported
checks are not counted as passes.}
\label{tab:22-quality-metrics}\centering\footnotesize
\setlength{\tabcolsep}{3pt}\renewcommand{\arraystretch}{1.12}
\begin{tabular}{@{}>{\raggedright\arraybackslash}p{.62\textwidth}ccc@{}}\toprule
Metric & Raw & Final & Controls\\\midrule
Valid molecules / all requests (\%) & \PendingCell & \PendingCell & \PendingCell \\
Connected molecules / all requests (\%) & \PendingCell & \PendingCell & \PendingCell \\
Exact L1 + all qualified design checks / requests (\%) & \PendingCell & \PendingCell & \PendingCell \\
Molecules with a justified chemical alert / requests (\%) & \PendingCell & \PendingCell & \PendingCell \\
Molecules requiring chemistry-context review / requests (\%) & \PendingCell & \PendingCell & \PendingCell \\
Required head retained / family-scope requests (\%) & \PendingCell & \PendingCell & \PendingCell \\
Correct role-local cycle allocation / requests (\%) & \PendingCell & \PendingCell & \PendingCell \\
Requested ring-size agreement / comparable requests (\%) & \PendingCell & \PendingCell & \PendingCell \\
TRAIN atom-environment coverage (\%) & \PendingCell & \PendingCell & \PendingCell \\
Ring-system TRAIN support / assessed products (\%) & \PendingCell & \PendingCell & \PendingCell \\
Exact L1 + product/component features observed / requests (\%) & \PendingCell & \PendingCell & \PendingCell \\
\bottomrule\end{tabular}
\par\smallskip\raggedright\scriptsize
TRAIN atom/ring support is a reference diagnostic and may share information with decoder priors.
Definitions: Table~\ref{tab:22-metric-definitions}; per-family outcomes:
Tables~\ref{tab:22-structure} and~\ref{tab:22-conditions}.
\end{table}

Agreement with lipid references is \PendingValue{realism outcome} in fingerprint and descriptor
spaces (Table~\ref{tab:22-realism-metrics}). Precision measures proximity to the reference
population, whereas coverage measures how much of that population is represented. Their joint
pattern is \PendingValue{precision-coverage interpretation}. Nearest-reference similarity and
internal diversity distinguish concentrated familiar outputs from diverse but distant structures.
The source-study-grouped two-sample classifier has AUC \PendingValue{two-sample AUC}, with
\PendingValue{two-sample uncertainty}. These comparisons distinguish TRAIN support from
independent lipid-reference agreement; near-chance discrimination does not prove chemical quality.

\begin{table}[H]
\caption{\textbf{Lipid realism and distributional coverage.}
Raw and final outputs are compared within each family and reference population, with
unique-molecule and all-request denominators distinguished. TRAIN-reference diagnostics are
separate from comparisons with source-disjoint observed lipids. Values are mean $\pm$ sample
standard deviation across training seeds.}
\label{tab:22-realism-metrics}\centering\footnotesize
\setlength{\tabcolsep}{3pt}\renewcommand{\arraystretch}{1.12}
\begin{tabular}{@{}>{\raggedright\arraybackslash}p{.49\textwidth}cc>{\raggedright\arraybackslash}p{.24\textwidth}@{}}\toprule
Metric & Raw & Final & Reference\\\midrule
Fingerprint precision (\%) & \PendingCell & \PendingCell & Declared lipid cohort \\
Fingerprint coverage (\%) & \PendingCell & \PendingCell & Declared lipid cohort \\
Fingerprint recall (\%) & \PendingCell & \PendingCell & Declared lipid cohort \\
Distinct reference-supported products / requests (\%) & \PendingCell & \PendingCell & Declared lipid cohort \\
Nearest-reference Tanimoto similarity & \PendingCell & \PendingCell & Declared lipid cohort \\
Internal ECFP4 diversity & \PendingCell & \PendingCell & Generated population \\
Descriptor distribution distance, per feature & \PendingCell & \PendingCell & Declared lipid cohort \\
Descriptor-space precision and coverage & \PendingCell & \PendingCell & Source-disjoint lipids \\
Source-grouped two-sample classifier AUC & \PendingCell & \PendingCell & Source-disjoint lipids \\
\bottomrule\end{tabular}
\par\smallskip\raggedright\scriptsize
ECFP4 uses Morgan radius 2 and 2,048 bits; neighborhood metrics use five neighbors.
Unsupported reference comparisons remain unassessed. Table~\ref{tab:22-descriptors} gives the
full descriptor panel; Table~\ref{tab:22-realism} reports each family and reference population.
\end{table}

\subsection{Recursive synthesis evidence and complete-route yield}
Exact final assembly closes L1 but does not establish an upstream route. Exact precursor-preparation
evidence (L2) covers \PendingValue{L2 component coverage} of required components, and qualified
terminal evidence (L3) covers \PendingValue{current L3 coverage}
(Table~\ref{tab:22-routes}). Complete-product dossier yield is
\PendingValue{complete dossier yield} per 1,000 requests, with
\PendingValue{family-wise evidence depth}. Each complete dossier requires every precursor branch
to terminate in an admitted exact material, preserving role, reaction-step and repeated-component
context. Directly accepted terminals do not require an L2 synthesis. Missing or scope-mismatched
sources, unresolved identities, expired availability and bounded-search censoring are reported
separately. The dominant limitation is \PendingValue{route-closure limitation}; absence of admitted
evidence is not evidence of unsynthesizability. These results concern auditable computational
synthesis dossiers and do not establish experimental synthesis or biological activity.


\section{Discussion}
FORGE incorporates synthesis constraints into whole-molecule generation without restricting the
neural generator to a fixed precursor catalogue. Across 22 reaction families, program conditioning
changes exact-assembly yield by \PendingValue{paired conditioned-minus-null contrast} percentage
points. Joint whole-molecule generation and component-factorized generation differ by
\PendingValue{family-wise joint-generation contrast}, indicating
\PendingValue{interpretation of coordination across precursor-derived regions}.

Leaving precursor identities open permits verified assemblies to contain components outside the
training catalogue. The observed component-novel yield of \PendingValue{yield}\% and held-component
recovery of \PendingValue{recovery}\% distinguish exploration beyond observed identities from
recovery of specifically withheld structures. Component novelty alone does not establish that a
product lies outside every catalogue-reachable assembly.

Exact assembly is one level of evidence. Chemical alerts, retained head sites, ring conditions,
reference-distribution coverage and component diversity assess different properties of the generated
lipids. Their joint results show \PendingValue{quality and diversity conclusion}, with residual
failures concentrated in \PendingValue{families and structural failure modes}. TRAIN-derived
structural priors contribute \PendingValue{inference gain and cost}; reference support computed
from those same priors cannot independently establish realism. Documentary L2/L3 assessment closes
\PendingValue{complete dossier yield}\% of requests and leaves
\PendingValue{unresolved and search-censored proportions}. These computational checks do not
establish experimental synthesis success, ionization or delivery performance.

% Main scientific content ends here; ICLR excludes the following disclosure statements.
\clearpage
\subsection*{AI use statement}
Generative AI tools assisted with conceptual and methodological development, software implementation
and debugging, experimental design and diagnostic analysis, interpretation of computational results,
literature search and summarization, and drafting and editing the manuscript. The authors are
responsible for the final code, reported results, citations and conclusions, including material
produced with AI assistance.

\subsection*{Ethics statement}
This work evaluates ionizable-lipid generation and synthesis evidence computationally. Molecular
design systems can be misused. Chemical validity, exact assembly, synthesis-route evidence and
predicted properties do not establish safety or biological efficacy. Generated structures require
appropriate evaluation before experimental use.

\subsection*{Reproducibility statement}
The Methods and Appendix specify the molecular representation, conditional discrete flow, numerical
sampler, exact assembly verification, route assessment, benchmark protocols and evaluation metrics.
Appendix~\ref{app:repro} records the source, configuration, split, checkpoint and evaluation artifacts.
Per-seed results appear in Table~\ref{tab:22-seeds}; assumptions and proofs appear in
Appendix~\ref{app:theory}.

\bibliographystyle{iclr2027_conference}
\bibliography{forge,references}
\clearpage
\appendix
\section{Method and implementation details}
\label{app:forge}

This appendix specifies program and graph encoding, model training, sampling, verification and
route assessment. Algorithm~\ref{alg:forge-inference} combines the inference steps.

\subsection{Reaction-program encoding}
\label{app:program-encoding}

The program tuple in Equation~\ref{eq:reaction-program} uses the ordered atom-position set
$\mathcal I_{\rm atom}(P)=\{1,\ldots,n(P)\}$. The arrays $\mathbf r$ and $\mathbf k$ assign
precursor roles and core positions to these atoms. This atom-position index differs from
$\mathcal I(P)$, which indexes all six families of categorical coordinates.

Morphology is computed from the selected atom-mapped trace and its sparse serialization.
For role $r$, let $V_r^{\rm ext}$ contain its exterior atoms, excluding reaction-core atoms, and
let $o_r(v)$ count the same-role exterior children of $v$ in the serialized tree. The four fields are
\begin{equation}
 M_r=(n_r,j_r,\gamma_r,a_r),\qquad
 n_r=|V_r^{\rm ext}|,\qquad
 j_r=\sum_{v\in V_r^{\rm ext}}\max\{o_r(v)-1,0\}.
\end{equation}
Here $\gamma_r$ counts closure edges with both endpoints in $V_r^{\rm ext}$, and $a_r$ counts
same-role bonds between exterior and core atoms. When the serialized tree restricts to a spanning
forest of the exterior subgraph, $\gamma_r$ equals that subgraph's cycle rank. The junction budget
is a tree-branching count; it does not count degree contributed by closure bonds. Consequently these
fields depend on the selected annotation and serialization, not on the unannotated graph alone.
Fixed-coordinate equality alone does not establish agreement with every morphology field.
Positions without precursor origin use the designated absent morphology state. Integer morphology
values are encoded with an offset of one, reserving zero for that absent state.

The training-fold program prior is defined in Section~\ref{sec:program-conditioning}.
The memory encoder uses separate program-identity and depth tokens,
$m_c=e_c(c)+e_{\rm type}(0)$ and $m_d=e_d(d)+e_{\rm type}(1)$, together with one token per active
atom position,
\begin{equation}
 m_i=e_r(r_i)+e_k(k_i)+e_{\rm pos}(i)+e_{\rm rep}(g_i)+e_{\rm comp}(u_i)
     +\sum_{v=1}^{4}e_{M,v}(M_{r_i,v})+e_{\rm type}(2).
\end{equation}
Here $g_i$ is the repeat-group state and $u_i$ the position within the precursor-origin component;
they are structural conditioning variables, not precursor identity labels. All embeddings share the
model hidden dimension. Layer normalization precedes memory self-attention, followed by residual
attention and feed-forward updates. The resulting program, depth and atom-position tokens provide
the memory for cross-attention in the graph denoiser. The morphology fields determine active
positions and fixed or variable masks. Sampling from the prior specifies these constraints;
generation assigns the atom, bond and topology states within each role.

A chemistry-specific adapter determines the precursor roles, reaction core, reverse decomposition
and forward transform associated with $c$. Its sampling layout determines active positions and
fixed or variable coordinates; the embeddings condition neural prediction on the remaining
program fields.

\subsection{Graph serialization and source annotations}
\label{app:graph-encoding}

A spanning-tree representation with at most twelve residual closure edges preserves the full
254-heavy-atom support without an $O(N_{\max}^2)$ edge state. The discrete flow generates the atom, bond
and topology states in this representation.
Let $\alpha=(\mathbf C,\tau,\mathbf r,\mathbf k)$ contain the selected exact precursor tuple,
registered assembly trace, atom-wise precursor-role assignment and reaction-core annotation. The deterministic
program-aware encoder
\begin{equation}
    \operatorname{Enc}_P:
    \{(G,\alpha):G\in\mathfrak G_{N_{\max},K_{\rm cl}}
      \text{ and }\alpha\text{ is admitted for }P\}
    \longrightarrow \mathfrak X_{N_{\max},K_{\rm cl}}
\end{equation}
constructs the categorical graph state
\begin{equation}
    X_\star:=\operatorname{Enc}_P(G,\alpha)=
    \left(
        \mathbf{a},
        \boldsymbol{\pi},
        \mathbf{b},
        \mathbf{e}^{L},
        \mathbf{e}^{R},
        \mathbf{z}
    \right)
\end{equation}
where $\mathbf{a}$ contains atom states, $\boldsymbol{\pi}$
parent pointers, $\mathbf{b}$ spanning-tree bond states, $\mathbf{e}^{L}$ and $\mathbf{e}^{R}$
closure-endpoint pointers, and $\mathbf{z}$ closure-bond states.

The encoder first partitions the atoms into connected precursor-origin components under $\alpha$.
It selects as root the reaction-core atom with the smallest tie-broken canonical RDKit rank, and
orders the origin-component graph by breadth-first traversal from that root component, breaking
ties by role index, component canonical ranks and original atom index. Within each component it uses
a deterministic depth-first preorder from the selected attachment atom, with neighbors ordered by
the same canonical ranks. The traversed bonds define the spanning tree. Every residual bond is
represented once as an ordered endpoint pair with its bond label, and residual triples are sorted
lexicographically. Role, core, fixed-coordinate and morphology arrays follow the
same atom order. Each admitted source record supplies a selected qualified exact assembly trace
and its precursor-origin assignment. Alternative source-role bindings remain explicit annotations;
they do not receive additional graph sampling mass when their product constitutions coincide. Thus
$\operatorname{Dec}(\operatorname{Enc}_P(G,\alpha))\cong G$ for every admitted record.

The raw padded categorical product space $\overline{\mathfrak X}_{N_{\max},K_{\rm cl}}$ also
contains malformed states. Its well-formed subset $\mathfrak X_{N_{\max},K_{\rm cl}}$ enforces an
active node prefix, earlier-parent pointers, distinct valid closure endpoints, no duplicate edges,
and null values in unused closure slots. The total decoder, including its failure value, is typed as
\begin{equation}
    \operatorname{Dec}:\overline{\mathfrak X}_{N_{\max},K_{\rm cl}}
    \longrightarrow
    \mathfrak G_{N_{\max},K_{\rm cl}}\cup\{\bot_{\rm dec}\}.
\end{equation}
For this study, $N_{\max}=254$ and $K_{\rm cl}=12$. The atom vocabulary includes qualified bromine states; bond states encode single, double and triple bonds, with aromatic structures represented through Kekul\'e bond assignments.
These bounds specify architectural support. Reaction consistency, nonzero probability under a
trained model and synthesis-program completeness require separate conditions.

The structural decoder specifies graph reconstruction from an already chosen state.
The production terminal readout additionally chooses that state from neural predictions
(Appendix~\ref{app:sampling}). These are separate operations.

The chemistry adapter marks any atom, pointer, or bond states fixed by $P$.
FORGE holds these states constant while corrupting the remaining states along the linear probability
path and predicting their clean values from $(X_t,t,P)$. Let $\Omega_P$ be the subset of the raw state
space satisfying the allowed coordinate alphabets and fixed values, and define the admissible program set
\begin{equation}
    \mathcal P_{\rm adm}
    =\left\{P:\mathfrak X_{N_{\max},K_{\rm cl}}\cap\Omega_P\neq\varnothing\right\}.
\end{equation}
The morphology prior is normalized only on $\mathcal P_{\rm adm}$. Each chemistry adapter is required
to satisfy the program-compatible encoding property: every admitted annotated record obeys
$\operatorname{Enc}_P(G,\alpha)\in\Omega_P$. Corpus admission requires this property, an exact
constitutional encode--decode round trip and source-consistent origin/core assignments. Records
failing admission are excluded. The frozen split is \PendingValue{split manifest}; noise marginals,
program priors and motif references are fitted only to its training partition. These checks do not
establish the property for unregistered chemistry.

\subsection{Transformer architecture and training objectives}
\label{app:training}

The denoiser is a graph Transformer over the complete sparse state. Each layer applies graph
self-attention with relation biases for parent, child and closure edges, then cross-attends to memory
tokens representing program identity, program depth, precursor role and reaction-core position.
Cross-attention is repeated in every layer so that reaction semantics can influence both local state
updates and long-range coordination among precursor-derived regions. The backbone has six layers,
hidden dimension 192, eight attention heads and dropout 0.1. Each routed adapter mixes three
experts with bottleneck dimension 64 using weights computed from the program summary. Additional
routed adapters feed atom/parent-bond predictions and closure predictions. Auxiliary heads predict
precursor roles, core positions and exterior child counts. The shared backbone receives
repeat-group, component-position and role-morphology conditioning.

Training assigns equal mass to the 22 admitted formal families and uniform mass to qualified
unique product constitutions within each family. Sampling is with replacement; multiple source
bindings do not multiply a constitution's probability, and sampling weights are not applied again
to the loss. Each optimizer update contains 3,168 examples from three distinct families, with
1,056 examples per selected family. A seeded permutation defines a balanced 22-update cycle in
which every family appears three times. Family batches are sharded across eight H100 GPUs, and
their coordinate reductions are accumulated before family-gradient projection. Padding is trimmed
to the longest molecule within a family batch while retaining the full closure support.

The optimizer is AdamW with learning rate $10^{-4}$, weight decay $5\times10^{-5}$,
$(\beta_1,\beta_2)=(0.9,0.999)$, $\epsilon=10^{-8}$ and global gradient-norm clipping at 1.
The update budget is \PendingValue{optimizer updates}; the independent training seeds and
checkpoint-selection rule are \PendingValue{training seeds and checkpoint rule}. Model and gradient
computations use deterministic FP32 arithmetic. Checkpoints retain model, optimizer, random-number
streams and the committed sampler state; uncommitted prefetch work is discarded on restart.
Training times are sampled uniformly, $t\sim\mathcal U(0,1)$, without endpoint clipping.

In the task losses below, $c$ indexes a formal family; multiple registered program variants in that
family share its task reduction.
For formal family $c$ and state family
$q\in\mathscr J=\{\mathbf a,\boldsymbol\pi,\mathbf b,\mathbf e^L,\mathbf e^R,\mathbf z\}$,
let $\mathcal S_{c,q}$ be the selected variable coordinates across that family's minibatch.
With coordinate loss $\ell_i=-\log d_{\theta,i}(X_{\star,i}\mid X_t,t,P)$, the pooled reduction is
\begin{equation}
 L_{c,q}^{\rm pool}=\frac{\sum_{i\in\mathcal S_{c,q}}\ell_i}{\max\{1,|\mathcal S_{c,q}|\}}.
\end{equation}
Coordinates here include their example index. Chemistry states (atoms, parent bonds and closure
bonds) additionally use equal mass over roles present in the selected coordinates. If
$\mathcal S_{c,q,r}$ contains coordinates assigned to role $r$ and
$\mathcal R_{c,q}=\{r:|\mathcal S_{c,q,r}|>0\}$, this reduction is
\begin{equation}
 L_{c,q}^{\rm role}=\frac{1}{|\mathcal R_{c,q}|}
 \sum_{r\in\mathcal R_{c,q}}\frac{\sum_{i\in\mathcal S_{c,q,r}}\ell_i}{|\mathcal S_{c,q,r}|},
 \label{eq:role-balanced-loss}
\end{equation}
with zero loss if no role is present. Atom and parent-bond groups use their position's role;
closure-bond groups use the role at the clean left endpoint. Pointer losses use the pooled reduction
and mask logits to the permitted layout positions. The flow term $L_{c,\rm flow}$ sums these six
reduced losses, with unit weights.

For each family, the task loss is
\begin{equation}
\begin{aligned}
 L_c={}&L_{c,\rm flow}
 +0.25L_{c,\rm role}+0.25L_{c,\rm core}+0.25L_{c,\rm repeat}\\
 &+L_{c,\rm child}+0.5L_{c,\rm junction}+L_{c,\rm chem\mid top}.
\end{aligned}
\label{eq:production-task-loss}
\end{equation}
The role and core terms are cross-entropies averaged equally over present target states.
The repeat term averages squared differences between predicted atom distributions, and separately
parent-bond distributions, over exact retained-fragment correspondences in repeated precursor
instances; it sums the two means. Atom correspondence requires a unique labelled fragment match
that preserves core annotations. Parent-bond correspondence matches both edge endpoints rather
than equating serialized child positions. Unequal fragments or ambiguous matches contribute no
repeat pairs; these target-derived alignments supervise the loss and are not supplied at inference.
The child term is state-balanced cross-entropy for same-component exterior child
counts in $\{0,1,2,3\}$. The junction term applies smooth-L1 loss to predicted and target
role-local junction totals, each divided by the number of exterior atoms in that role, then averages
over eligible examples within a role and equally across present roles. Predicted junction totals
sum the expected child-count contribution $\max\{o-1,0\}$ across exterior atoms.
The topology-conditioned chemistry term repeats the role-balanced atom and bond losses in a second
forward pass with clean parent and closure endpoints and corrupted atom and bond states.
Each ablation changes only the declared objective or architecture terms under its matched training contract.

PCGrad combines the task gradients in integer formal-family-state order. For
$g_c=\nabla_\theta L_c$, it initializes $\widetilde g_c=g_c$ and, for each other task $h$ in that
order, applies
\begin{equation}
 \widetilde g_c\leftarrow\widetilde g_c-
 \mathbf1\{\langle\widetilde g_c,g_h\rangle<0\}
 \frac{\langle\widetilde g_c,g_h\rangle}{\max\{\|g_h\|_2^2,10^{-12}\}}\,g_h
\end{equation}
on parameter coordinates available to both tasks. Gradients are not norm-normalized, and the
optimizer receives the arithmetic mean of the projected task gradients. This update is not generally
the gradient of a scalar sum of task losses. No population optimality claim is made for the full
procedure. In particular, target-dependent group reductions and the auxiliary objectives require
analysis beyond Proposition~\ref{prop:masked-bayes}.

Layerwise cross-attention, auxiliary losses, routed adapters and gradient-conflict control are
separate interventions in the architecture comparison (Table~\ref{tab:22-architecture}).

\subsection{Numerical sampling and terminal decoding}
\label{app:sampling}

The sampler uses \PendingValue{flow steps} steps, denoted $L_{\rm samp}$. Its strictly positive atom and bond
source marginals are indexed by registered program and precursor role and fitted to TRAIN data.
Pointer marginals are restricted and normalized on their allowed layout positions. The finite
sampler evaluates the denoiser at $t_k=k/L_{\rm samp}$, for
$k=0,\ldots,L_{\rm samp}-1$, and uses $\Delta t=1/L_{\rm samp}$.
Conditional on $Z_k=x$, it draws one clean category
$A_{k,i}\sim d_{\theta,i}(\cdot\mid x,t_k,P)$ independently for each variable coordinate.
Write $s_{t,i}^P(u\mid a)=(1-t)q_{0,i}^P(u)+t\mathbf1\{u=a\}$ for its conditional path factor,
$\dot s_i^P(u\mid a)=\mathbf1\{u=a\}-q_{0,i}^P(u)$ for the derivative, and
$[z]_+=\max\{z,0\}$. The code uses the conditional off-diagonal rate
\begin{equation}
 R_{t,i}^{*,\varepsilon}(u,v\mid a,P)
 =\frac{[\dot s_i^P(v\mid a)-\dot s_i^P(u\mid a)]_+}
 {|\mathcal X_i(P)|\max\{s_{t,i}^P(u\mid a),\varepsilon\}},
 \quad u\neq v,\qquad \varepsilon=10^{-8}.
 \label{eq:clamped-rate}
\end{equation}
Full source support is required only on the allowed coordinate alphabet, including the masked
candidate set for a pointer. Thus that alphabet is the support for every $t_k<1$.
The denominator clamp in Equation~\ref{eq:clamped-rate} is a numerical safeguard absent from the
ideal rate in Appendix~\ref{app:flow-proof}.

For the sampled endpoint $a=A_{k,i}$, let
$h_{k,i}=\Delta t\sum_{v\neq x_i}R_{t_k,i}^{*,\varepsilon}(x_i,v\mid a,P)$ and set
$\kappa_{k,i}=\min\{1,0.999/h_{k,i}\}$ when $h_{k,i}>0$, with $\kappa_{k,i}=1$ otherwise.
The implemented coordinate kernel is
\begin{equation}
 \widehat K_{k,i}(v\mid x,a)=
 \begin{cases}
 \kappa_{k,i}\Delta t\,R_{t_k,i}^{*,\varepsilon}(x_i,v\mid a,P),&v\neq x_i,\\
 1-\kappa_{k,i}h_{k,i},&v=x_i.
 \end{cases}
 \label{eq:capped-kernel}
\end{equation}
Its off-diagonal masses sum to at most $0.999$; the diagonal is the remaining mass. Consequently
all entries are nonnegative and each row sums to one. Coordinates are drawn independently from
these kernels, followed by the projection $\Pi_P$.

Without the clamp or cap, and when every conditional Euler row is nonnegative, averaging over
$A_{k,i}$ gives the Euler kernel built from Equation~\ref{eq:posterior-rate} exactly. One sampled
endpoint per coordinate therefore does not by itself bias that uncapped one-step law. The nonlinear
cap generally changes the averaged kernel. Independent coordinate updates also permit simultaneous
finite-step changes, so their product is a numerical approximation to the local CTMC, not its exact
transition kernel. No discretization error bound is asserted near the terminal singularity.

After all steps, the model evaluates $(Z_{L_{\rm samp}},1,P)$ once more. The production terminal
readout $H_P^\theta$ selects coordinate values using constrained argmax, restores fixed states and
returns either a raw state or a typed terminal failure, denoted collectively by $\bot_{\rm term}$.
The atom-aware readout first chooses atom states subject to unavoidable incident-bond demands
from the fixed core and non-root tree attachments; subsequent variable bonds must fit those atom
capacities. It abstains when this bounded assignment cannot be completed. It does not exchange an
element for a higher-valence element merely to accommodate a chosen topology. The registered
core condition reserves the required bond and hydrogen states for each program. In the AGILE-type
Ugi 3-CR, these include one hydrogen on the aldehyde-derived $\alpha$ carbon and one on the
isocyanide-derived amide nitrogen; these conditions are not imposed on other Ugi variants.
The selected state is then reconstructed and checked for
connectedness, valence and sanitization. A readout failure retains its detailed reason in the ledger;
for graph-law notation we extend $\operatorname{Dec}(\bot_{\rm term})=\bot_{\rm dec}$.
If $q_{\theta,\rm out}^P$ is the raw law after terminal readout, including failure, then
\begin{equation}
 p_\theta(\cdot\mid P)=\operatorname{Dec}_{\#}q_{\theta,\rm out}^P,
 \qquad
 p_\theta(G\mid P)=\sum_{z:\operatorname{Dec}(z)=G}q_{\theta,\rm out}^P(z).
 \label{eq:generated-law}
\end{equation}
This law retains structural failures and chemically invalid graphs before the validity and L1
filters; it is not conditioned on acceptance.
The ideal theorem does not include the learned $t=1$ readout, which is unconstrained by a loss
integrated over $t\in(0,1)$. Decoder and source-prior comparisons use their declared fixed-weight
or retrained contrasts (Table~\ref{tab:22-decoder}).

\subsubsection{Checked construction and selection}

Let $\Pi$ specify a bounded inference policy: the draw count, ordered proposal operators, TRAIN
references, search limits, source verifier, selection objective and selection-group size. Its
version is \PendingValue{inference policy}; the draw count $D$ is \PendingValue{draws per request},
the stage caps are \PendingValue{proposal and search caps}, and the selection-group size is
\PendingValue{selection batch size}. For request $b$, retain all raw trajectories and form
\begin{equation}
 \mathcal C_b=\{G^{\rm raw}_{b,d}:1\leq d\leq D\}
 \cup\mathcal R_\Pi\!\left(P_b,\{Z^{\rm out}_{b,d},\lambda_{b,d}\}_{d=1}^{D}\right),
\end{equation}
where $\lambda_{b,d}$ denotes saved terminal predictions and $\mathcal R_\Pi$ is the bounded
construction map. Original graphs, failed trajectories, rejected proposals and exceptions remain
in the attempt ledger. A request is not removed when its candidate set has no eligible graph.

Construction can tie repeated components, restore source-required component constraints, propose
qualified retained-head sites, and assign TRAIN-supported joint atom/bond or ring-system motifs.
The operators act on generated atom slots while preserving the declared atom budget, precursor
origins and fixed reaction core. Atom relabelling and topology edits are recorded separately.
Ring repair is applied after component and head construction, and all checks are evaluated again
on the final graph. Repeated arms receive tied edits when their established correspondences permit
this; every proposed product must still pass the unchanged source verifier. A TRAIN motif is an
explicit inference prior and does not itself certify chemical plausibility or a complete route.
An unsupported or context-dependent ring fragment is unassessed rather than assigned a fabricated
known identity. Graph cycle rank, ring systems, SSSR rings and serialization-dependent fundamental
cycles remain distinct quantities.

The selector maps the request-indexed pools to at most one graph per request,
\begin{equation}
 (G^{\rm sel}_1,\ldots,G^{\rm sel}_B)
 =\mathcal S_\Pi\bigl((P_b,\mathcal C_b)_{b=1}^{B}\bigr),
 \qquad G^{\rm sel}_b\in\mathcal C_b\cup\{\bot_{\rm dec}\}.
\end{equation}
Source-exact eligibility, applicable design predicates and product/role-component diversity and
novelty constraints are explicit parts of $\Pi$. A combined design pass requires every applicable
qualified predicate to be assessed and pass; unassessed checks cannot establish that endpoint.
The selection objective and constraint reference are \PendingValue{selector objective and floors}.
Each output retains its request, draw and proposal lineage. Raw and selected yields both use the
original $B$ requests as denominator, while cost accounting includes $BD$ trajectories, all
construction attempts, source checks and selection work.

The selected joint law is the pushforward of the raw trajectory law through
$\mathcal R_\Pi$ and $\mathcal S_\Pi$. It is distinct from Equation~\ref{eq:generated-law} and
can couple different requests. The ideal flow theorem does not certify these inference maps;
independent-draw accepted-count formulas apply only when their independence assumptions hold.
A graph-changing repair changes the proposal distribution, whereas rejecting completed graphs
alone does not.

\subsection{Exact assembly verification}
\label{app:verification}

For the chemistry $c$ specified by $P$, FORGE decomposes each valid graph into role-assigned
precursors and applies the registered forward transform. Exact replay checks assembly existence
under $c$; it does not test full agreement with the sampled conditioning layout.

\begin{definition}[Exact assembly and recovered witnesses]
\label{def:exact-assembly}
For the registered forward operator $\mathcal F_c$ on precursor inputs $\mathfrak C_c$, define
\begin{equation}
 \mathcal E_c^\star(G)=\{\mathbf C\in\mathfrak C_c:G\in\mathcal F_c(\mathbf C)\},\qquad
 \mathrm{L1}^\star(G,c)=\mathbf1\{\mathcal E_c^\star(G)\neq\varnothing\}.
\end{equation}
The bounded reverse procedure returns candidate pairs $(\mathbf C,\tau)$ in $\mathcal D_c(G)$.
Equation~\ref{eq:main-exact-traces} retains the pairs passing exact forward replay, and
\begin{equation}
 \widehat{\mathrm{L1}}(G,c)=\mathbf1\{\widehat{\mathcal E}_c(G)\neq\varnothing\}.
\end{equation}
Both predicates describe assembly existence; molecular validity is checked separately.
\end{definition}
The set $\mathcal E_c^\star(G)$ contains precursor inputs, whereas
$\widehat{\mathcal E}_c(G)$ contains input--trace pairs. Their compatible comparison is
\begin{equation}
 \{\mathbf C:\exists\tau,\ (\mathbf C,\tau)\in\widehat{\mathcal E}_c(G)\}
 \subseteq\mathcal E_c^\star(G),\qquad
 \widehat{\mathrm{L1}}(G,c)\leq\mathrm{L1}^\star(G,c).
\end{equation}
The inclusion follows directly from exact replay. Its converse requires complete reverse search,
which has not been established. Throughout the paper, verified exact-L1 yield counts valid
attempts satisfying the recovered predicate. Candidate traces are compared after the registered
role- and step-aware canonicalization, so duplicate encodings do not create additional witnesses.

Post-hoc filtering changes which completed samples are retained but cannot change the proposal
distribution's exact-L1 acceptance probability.
Under independent sampling, this probability determines the complete distribution of accepted
counts under a fixed attempt budget and the expected cost of obtaining a fixed number of accepted
products (Proposition~\ref{prop:posthoc-budget}).
We distinguish decomposition coverage from forward-replay precision because a candidate disconnection
may fail to reconstruct $G$.
Candidate decomposition coverage is the fraction of valid graphs for which
$\mathcal D_c(G)\neq\varnothing$. Candidate yield additionally uses all original requests as
its denominator; undefined conditional ratios are not replaced by zero. For a declared sequence of valid evaluation graphs $G_1,\ldots,G_m$, candidate replay precision is
\begin{equation}
 \mathrm{ReplayPrecision}
 =\frac{\sum_{b=1}^m|\widehat{\mathcal E}_c(G_b)|}
 {\sum_{b=1}^m|\mathcal D_c(G_b)|},
\end{equation}
with N/E when the denominator is zero. Replay of tuples already in
$\widehat{\mathcal E}_c(G)$ has $100\%$ success by construction. This invariant does not estimate
candidate replay precision. We separately define
\begin{equation}
 \mathrm{CandidateAmbig}(G,c)=\mathbf 1\{|\mathcal D_c(G)|>1\},\qquad
 \mathrm{ExactAmbig}(G,c)=\mathbf 1\{|\widehat{\mathcal E}_c(G)|>1\}.
\end{equation}
A distinct exact-L1 product is a deduplicated generated constitution. An unambiguous exact
decomposition satisfies $|\widehat{\mathcal E}_c(G)|=1$. We distinguish product uniqueness from
decomposition uniqueness throughout.

\subsection{Recursive route assessment and dossier admission}
\label{app:route-assessment}

\subsubsection{Assessment interface}

For a valid generated graph $G$ and reaction program $P$ with chemistry $c$, route assessment
returns either a
complete synthesis dossier or a typed abstention.
Let $\mathfrak B$ denote the frozen route contract: depth and query budgets, registered transforms,
evidence-admission rules and a dated supplier snapshot. Let $\mathfrak D$ be the space of candidate
dossiers, each containing a synthesis tree and its reaction and terminal evidence, and let
$\mathfrak F_{\mathfrak B}$ be the disjoint finite set of typed failures.
Generation and assessment are defined by
\begin{equation}
    Y=\mathsf{Assess}_{\mathfrak B}(G,P,\mathbf C_T,\tau_T)
    \in\mathfrak D\cup\mathfrak F_{\mathfrak B}.
\end{equation}
$\mathsf{Assess}_{\mathfrak B}$ is invoked after the verifier returns one admitted root trace
$(\mathbf C_T,\tau_T)\in\widehat{\mathcal E}_c(G)$. It verifies the final assembly, searches for
preparation routes to unavailable precursors and checks terminal supplier evidence. $P$ and the assessor act at different stages: $P$ guides generation, whereas the assessor
constructs a dossier from the completed graph without modifying it.
FORGE accepts $(G,T)$ only when $G$ passes molecular and exact-assembly verification and $T$ traces
every required precursor to an available starting material; otherwise, it abstains.

\subsubsection{Recursive expansion and completion checks}

Completing a synthesis dossier requires evidence for every recovered precursor. Exact L1
verification specifies the precursor structures; availability and preparation routes require
additional evidence. Each precursor must therefore match an available material or have a supported
route from available starting materials.
Let $T$ be a candidate dossier with bipartite synthesis tree $(V_M,V_R,E_T)$ rooted at $G$, where each product molecule
points to its preparation reaction and each reaction points to its required reactants.
The assessor constructs these routes recursively and withholds a complete dossier when any branch
remains unresolved.
A molecule node terminates when a current, dated supplier record lists an exact constitutional match
as available; otherwise, FORGE must expand that node through a supported, forward-verified reaction.
FORGE uses a bounded level-two (L2) planner to propose each expansion.
For each proposal, FORGE records the reaction, reactants, product, and supporting evidence, and
accepts the step only when the forward transform reconstructs the target precursor exactly.
Each accepted reaction adds its reactants as new molecule nodes, and FORGE applies the same
termination-or-expansion rule recursively.
The current assessor admits only an unambiguous verified exact-L1 root. Every dossier therefore
records the sole selected root trace
$(\mathbf C_T,\tau_T)\in\widehat{\mathcal E}_c(G)$; an ambiguous exact set produces a typed
abstention. For each upstream reaction node $r$, let
$\mathrm{Admitted}_{\mathfrak B}(r)=1$ when its reaction and substrate scope satisfy the frozen
evidence policy, and let $\mathrm{Verify}_{\mathfrak B}(r)=1$ when
$\operatorname{target}(r)\in\mathcal F_r(\operatorname{reactants}(r))$.
For each leaf molecule $\ell$, let $\mathrm{Available}_{\mathfrak B}(\ell)=1$ when it satisfies the terminal
supplier-evidence rule.
Let $V_R^{\mathrm{up}}(T)$ denote the upstream preparation-reaction nodes and $L(T)$ the leaf
molecules. $\mathrm{WellFormed}(G,T,P)$ requires that $T$ is a connected, acyclic bipartite tree
rooted at $G$; each internal molecule has exactly one recorded producing reaction; no molecular constitution repeats along an ancestor path; each reaction's
molecule children equal its recorded reactants with their multiplicities; and $T$ has no disconnected, cyclic or
extraneous node. Its distinguished final-assembly record must have product $G$ and precursor
children exactly $\mathbf C_T$ under trace $\tau_T$.
\begin{equation}
\begin{aligned}
    \mathrm{Complete}_{\mathfrak B}(G,T,P)
    ={}&\mathrm{Valid}(G)\,\mathrm{WellFormed}(G,T,P)\\
    &\times\mathbf1\!\left[\widehat{\mathcal E}_c(G)=\{(\mathbf C_T,\tau_T)\}\right]\\
    &\times\prod_{r\in V_R^{\mathrm{up}}(T)}
    \mathrm{Admitted}_{\mathfrak B}(r)\mathrm{Verify}_{\mathfrak B}(r)
    \prod_{\ell\in L(T)}\mathrm{Available}_{\mathfrak B}(\ell).
\end{aligned}
\end{equation}
All factors are binary and empty products equal one. Repeated use of a material is represented by
separate tree occurrences; equality of their constitutions does not merge their nodes.
Under finite search bounds and terminating evidence queries, the assessor terminates and returns a
dossier only when $\mathrm{Complete}_{\mathfrak B}(G,T,P)=1$
(Lemma~\ref{lem:planner-termination} and Invariant~\ref{inv:dossier-soundness}).
An abstention means that FORGE did not complete a dossier under $\mathfrak B$; it does not imply that
no synthesis route exists outside the frozen search space or evidence snapshot.
The guarantee is computational and does not establish experimental synthesis success.

\subsubsection{Recorded evidence and support tiers}

For each accepted pair $(G,T)$, the synthesis dossier contains every molecule and reaction in $T$,
the evidence for each reaction and a dated supplier record for each terminal material.
We classify assembly support relative to the fixed reference space into four tiers:
\textbf{E0 (enumerated):} the generated lipid matches a product that was already constructed
computationally from known components using a supported reaction.
\textbf{E1 (known components):} the generated lipid was not included in the bounded E0 enumeration,
but a supported reaction assembles it entirely from known components.
\textbf{E2 (generated precursor):} a supported reaction assembles the lipid, but at least one
precursor is newly generated and requires a complete L2 route to available starting materials.
\textbf{E3 (new assembly family):} the lipid requires a final assembly reaction outside the
qualified reaction registry.

\subsection{End-to-end inference procedure}
\label{app:inference}

Algorithm~\ref{alg:forge-inference} records every requested output and every trajectory used to
construct it, including failures in molecular validation, exact replay or route completion.

\begingroup
\setlength{\parskip}{0pt}
\hrule height 0.8pt\kern 2pt
\refstepcounter{algorithm}
\noindent\textbf{Algorithm~\thealgorithm} FORGE generation and fail-closed synthesis-program construction.
\label{alg:forge-inference}
\par\nobreak\kern 2pt\hrule height 0.4pt\nobreak
\begin{algorithmic}[1]
\Require trained coordinate predictors, chemistry $c$, program prior $p_{\rm prog}(\cdot\mid c)$,
integers $L_{\rm samp},D\geq1$ and $B\geq0$, recorded seed $s$, inference policy $\Pi$,
route contract $\mathfrak B$
\Ensure attempt-indexed dossier sequence $\mathcal D_{\rm out}$, exact-L1 ledger $\mathcal L_{\rm out}$,
and typed failure and proposal ledger $\mathcal H$
\State initialize the random-number streams from $s$;
$\mathcal D_{\rm out}\gets[\,]$; $\mathcal L_{\rm out}\gets[\,]$; $\mathcal H\gets[\,]$
\For{$b=1,\ldots,B$}
    \State sample and record $P_b\sim p_{\rm prog}(\cdot\mid c)$; initialize raw draw ledger $\mathcal Z_b$
    \For{$d=1,\ldots,D$}
        \State sample $Z_0\sim q_0^{P_b}$ and set $Z\gets\Pi_{P_b}(Z_0)$
        \For{$k=0,\ldots,L_{\rm samp}-1$}
            \State $t_k\gets k/L_{\rm samp}$; $\Delta t_k\gets1/L_{\rm samp}$
            \State $Z\gets\Pi_{P_b}\!\left(\Call{CappedCoordinateEuler}{Z,t_k,\Delta t_k,P_b}\right)$
        \EndFor
        \State obtain terminal predictions $\lambda_{b,d}$ and $Z^{\rm out}_{b,d}\gets H_{P_b}^\theta(Z)$
        \State append the state or typed terminal failure, predictions and raw decoded graph to $\mathcal Z_b$
    \EndFor
    \State $\mathcal C_b\gets\Call{CheckedConstruction}{P_b,\mathcal Z_b,\Pi}$
    \State record every raw/proposed graph, check, abstention and cost in $\mathcal H$
\EndFor
\State $(G_1,\ldots,G_B)\gets\mathcal S_\Pi((P_b,\mathcal C_b)_{b=1}^{B})$
\For{$b=1,\ldots,B$}
    \State $G\gets G_b$; $P\gets P_b$
    \If{$G=\bot_{\rm dec}$ or $\mathrm{Valid}(G)=0$}
        \State record $(b,\textsc{InvalidGraph},\text{decode or validation reason})$ in $\mathcal H$
        \State \textbf{continue}
    \EndIf
    \State $\widehat{\mathcal{E}}_c(G)\gets
    \{(\mathbf C,\tau)\in\mathcal D_c(G):\exists H\in\mathcal F_c(\mathbf C),\ H\cong G\}$
    \If{$\widehat{\mathcal{E}}_c(G)=\varnothing$}
        \State $\mathcal H\gets\mathcal H\mathbin{\|}[(b,\textsc{L1Failure})]$
        \State \textbf{continue}
    \EndIf
    \State $\mathcal L_{\rm out}\gets\mathcal L_{\rm out}\mathbin{\|}
    [(b,G,\widehat{\mathcal E}_c(G))]$
    \If{$|\widehat{\mathcal E}_c(G)|\neq1$}
        \State $\mathcal H\gets\mathcal H\mathbin{\|}[(b,\textsc{AmbiguousL1})]$
        \State \textbf{continue}
    \EndIf
    \State select the sole $(\mathbf C_T,\tau_T)\in\widehat{\mathcal E}_c(G)$
    \State $Y\gets\mathsf{Assess}_{\mathfrak B}(G,P,\mathbf C_T,\tau_T)$
    \If{$Y$ is a dossier $T$ and $\mathrm{Complete}_{\mathfrak B}(G,T,P)=1$}
        \State $\mathcal D_{\rm out}\gets\mathcal D_{\rm out}\mathbin{\|}[(b,T)]$
    \ElsIf{$Y\in\mathfrak F_{\mathfrak B}$}
        \State $\mathcal H\gets\mathcal H\mathbin{\|}[(b,Y)]$
    \Else
        \State $\mathcal H\gets\mathcal H\mathbin{\|}[(b,\textsc{IncompleteDossier})]$
    \EndIf
\EndFor
\State \Return $\mathcal D_{\rm out},\mathcal L_{\rm out},\mathcal H$
\end{algorithmic}
\nopagebreak\hrule height 0.8pt
\endgroup

\subsection{Reproducibility and artifact provenance}
\label{app:repro}

Each run is identified by its training seed, data and split manifests, atom/program vocabularies,
source registries, fitted noise marginals, configuration, code snapshot and checkpoint SHA-256.
The independent training-seed identifiers are \PendingValue{training seed identifiers}. Each
comparison retains its declared attempt budget, checkpoint rule and evaluator. Sampling seeds,
request/layout identifiers and proposal lineage are recorded separately from training seeds.
The manuscript provenance records identify the source artifacts for tables and figures through
SHA-256 hashes.

The inference-policy identifier is \PendingValue{inference policy identifier}; its record includes
all proposal caps, TRAIN motif references, admission order, selection-group boundaries and source
checks. The route contract is \PendingValue{route contract identifier}, including depth/query
limits, evidence-admission policy and dated terminal-material snapshots. Restarts preserve committed
checkpoints and completed request shards, and source-preserving fixes retain the exact earlier
implementation and an equivalence receipt. A change to an inference map defines a distinct arm
unless equivalence is demonstrated for the declared inputs.

Reproduction requires source code, configuration and environment manifests, model checkpoints,
split/registry hashes, attempt-level output and verification ledgers, and figure/table exporters.
These artifacts are identified by \PendingValue{reproduction archive identifier}. Corpus access conditions are
\PendingValue{data availability}. Constitutional identity removes atom maps and stereochemical
labels while retaining formal charge and bond order; no tautomer normalization is applied.

\clearpage


\clearpage
\section{Evaluation across reaction families}
\label{app:22-results}
We evaluate complete lipid graphs at three distinct levels: exact assembly consistency, structural
and distributional properties, and synthesis-evidence completeness. Unless stated otherwise, $N$
denotes all requested products. Conditional rates specify their assessed population and retain
invalid, unsupported and abstained requests in the accompanying accounting.

\subsection{Cohort, program registry and split accounting}
We associate each registered assembly program with its source evidence, reactive-site multiplicity,
precursor roles and repeat/depth support. Source rows and unique constitutional products are counted
separately. The family census includes sampling mass, component inventories, source-study overlap
and exclusions; support tiers E0--E3 describe evidence depth rather than a probability of synthesis
success.
\begin{table}[H]
\caption{\textbf{Cohort census.} Source rows and constitutional products are separate units. Source and role inventories report sampling mass, component counts, source-study overlap and exclusions.}
\label{tab:22-census}
\centering\footnotesize\setlength{\tabcolsep}{3pt}\renewcommand{\arraystretch}{1.12}
\begin{tabular}{@{}>{\raggedright\arraybackslash}p{.30\textwidth}ccccc@{}}
\toprule
Family & \shortstack{Admitted\\records} & \shortstack{Unique\\products} & \shortstack{Programs /\\studies} & \shortstack{TRAIN /\\CAL / TEST} & \shortstack{Excluded /\\unassessed} \\
\midrule
A3 amine--aldehyde--alkyne & \PendingCell & \PendingCell & \PendingCell & \PendingCell & \PendingCell \\
Acid--epoxide diester & \PendingCell & \PendingCell & \PendingCell & \PendingCell & \PendingCell \\
AEMA & \PendingCell & \PendingCell & \PendingCell & \PendingCell & \PendingCell \\
Aldehyde Ugi 3-CR & \PendingCell & \PendingCell & \PendingCell & \PendingCell & \PendingCell \\
Aldehyde Ugi 4-CR & \PendingCell & \PendingCell & \PendingCell & \PendingCell & \PendingCell \\
$\alpha$-Isocyanoester dihydroimidazole & \PendingCell & \PendingCell & \PendingCell & \PendingCell & \PendingCell \\
Amine alkylation & \PendingCell & \PendingCell & \PendingCell & \PendingCell & \PendingCell \\
Amine--epoxide & \PendingCell & \PendingCell & \PendingCell & \PendingCell & \PendingCell \\
Aryl reductive amination & \PendingCell & \PendingCell & \PendingCell & \PendingCell & \PendingCell \\
Aza-Michael acrylamide & \PendingCell & \PendingCell & \PendingCell & \PendingCell & \PendingCell \\
Aza-Michael acrylate & \PendingCell & \PendingCell & \PendingCell & \PendingCell & \PendingCell \\
Disulfide Michael & \PendingCell & \PendingCell & \PendingCell & \PendingCell & \PendingCell \\
Epoxide O-acylation & \PendingCell & \PendingCell & \PendingCell & \PendingCell & \PendingCell \\
iPhos & \PendingCell & \PendingCell & \PendingCell & \PendingCell & \PendingCell \\
Ketone--isocyanide amide & \PendingCell & \PendingCell & \PendingCell & \PendingCell & \PendingCell \\
Ketone Ugi 4-CR & \PendingCell & \PendingCell & \PendingCell & \PendingCell & \PendingCell \\
Maleate & \PendingCell & \PendingCell & \PendingCell & \PendingCell & \PendingCell \\
O-esterification & \PendingCell & \PendingCell & \PendingCell & \PendingCell & \PendingCell \\
Passerini & \PendingCell & \PendingCell & \PendingCell & \PendingCell & \PendingCell \\
Preassembled thiol-yne & \PendingCell & \PendingCell & \PendingCell & \PendingCell & \PendingCell \\
Reductive amination & \PendingCell & \PendingCell & \PendingCell & \PendingCell & \PendingCell \\
Thiolactone & \PendingCell & \PendingCell & \PendingCell & \PendingCell & \PendingCell \\
\bottomrule
\end{tabular}
\end{table}

\subsection{Training, numerical validation and computational cost}
\begin{table}[H]
\caption{\textbf{Training and evaluation settings.} Settings and run identifiers are reported separately for each independent training seed.}
\label{tab:22-training}
\centering\small
\begin{tabular}{@{}p{.70\textwidth}l@{}}\toprule
Field & Setting / result\\\midrule
Cohort, split, source registry and vocabulary hashes & \PendingCell\\
Checkpoint IDs, independent seed IDs and stopping rule & \PendingCell\\
Per-family sampling law, presentations and update appearances & \PendingCell\\
Global/per-device batch, accumulation and optimizer updates & \PendingCell\\
Optimizer, learning-rate schedule, loss weights and clipping & \PendingCell\\
Architecture, adapters, role/core/repeat and topology objectives & \PendingCell\\
Precision, determinism, devices, workers and memory & \PendingCell\\
Numerical equivalence, kill/restart and checkpoint provenance & \PendingCell\\
Training, setup, checkpoint, sampling and verifier time/cost & \PendingCell\\
Requests per family/seed, saved draw IDs and final inference policy & \PendingCell\\
Sampled-layout distribution and all method-visible information & \PendingCell\\
\bottomrule\end{tabular}
\end{table}
Training diagnostics resolve atom, bond, parent, closure and repeated-component errors by family
and corruption time. Checkpoint comparisons use fixed development requests and report exact yield,
structural assessment and diversity alongside gradients and clipping. We distinguish molecular
presentations, family appearances in updates and optimizer updates. Matched-exposure batch/update
comparisons and exposure-scaling comparisons isolate different effects. Checkpoint selection uses
\PendingValue{checkpoint-selection rule}.
\begin{figure}[H]\centering
\ResultPanel{a. Learning and exposure}{Per-family losses and errors against presentations and optimizer updates.}\hfill
\ResultPanel{b. Quality and cost}{Fixed development checkpoint metrics; latency/calls per accepted distinct product.}
\caption{\textbf{Training adequacy and computational efficiency.} Learning curves distinguish presentations from optimizer updates. Hardware, precision, deterministic settings, numerical failures and elapsed budget are reported for each arm.}
\label{fig:22-training}\end{figure}

\subsection{Exact-L1 decomposition and failure accounting}
We report exact assembly outcomes for each family, registered program, training seed and inference
arm. Each request contributes to the denominator whether generation, decomposition or verification
succeeds. Alongside exact yield, we count valid and connected graphs, distinct exact products,
unique exact precursor roots and checked decomposition traces.
Candidate-search coverage is the fraction of requests for which a decomposition candidate is found;
verified decomposition coverage requires exact replay. Candidate precision counts passing checked
traces over all checked traces, with ambiguity and search censoring explicit. The replay rate among
already accepted exact traces is an implementation invariant, not an independent precision estimate.
We distinguish family-local product counts from counts deduplicated across all families. Equal-family
macro averages, pooled yields and worst-family outcomes use separate definitions. Training-seed
variability accompanies the family-wise estimates.

\begin{table}[H]
\caption{\textbf{Attempt-level outcomes.} Raw and selected outputs are evaluated separately for each registered program and seed. Malformed graphs, support violations, absent decompositions, rejected source domains, exceptions and search censoring are separate outcomes. Counts retain all-request denominators; nonexclusive events may overlap.}
\label{tab:22-outcomes}
\centering\footnotesize\setlength{\tabcolsep}{3pt}\renewcommand{\arraystretch}{1.12}
\begin{tabular}{@{}>{\raggedright\arraybackslash}p{.30\textwidth}ccccc@{}}\toprule
Family & \shortstack{Valid /\\requests} & \shortstack{Connected /\\requests} & \shortstack{Distinct\\exact / $N$} & \shortstack{Ambiguous\\exact roots} & \shortstack{Failure /\\abstention} \\\midrule
A3 amine--aldehyde--alkyne & \PendingCell & \PendingCell & \PendingCell & \PendingCell & \PendingCell \\
Acid--epoxide diester & \PendingCell & \PendingCell & \PendingCell & \PendingCell & \PendingCell \\
AEMA & \PendingCell & \PendingCell & \PendingCell & \PendingCell & \PendingCell \\
Aldehyde Ugi 3-CR & \PendingCell & \PendingCell & \PendingCell & \PendingCell & \PendingCell \\
Aldehyde Ugi 4-CR & \PendingCell & \PendingCell & \PendingCell & \PendingCell & \PendingCell \\
$\alpha$-Isocyanoester dihydroimidazole & \PendingCell & \PendingCell & \PendingCell & \PendingCell & \PendingCell \\
Amine alkylation & \PendingCell & \PendingCell & \PendingCell & \PendingCell & \PendingCell \\
Amine--epoxide & \PendingCell & \PendingCell & \PendingCell & \PendingCell & \PendingCell \\
Aryl reductive amination & \PendingCell & \PendingCell & \PendingCell & \PendingCell & \PendingCell \\
Aza-Michael acrylamide & \PendingCell & \PendingCell & \PendingCell & \PendingCell & \PendingCell \\
Aza-Michael acrylate & \PendingCell & \PendingCell & \PendingCell & \PendingCell & \PendingCell \\
Disulfide Michael & \PendingCell & \PendingCell & \PendingCell & \PendingCell & \PendingCell \\
Epoxide O-acylation & \PendingCell & \PendingCell & \PendingCell & \PendingCell & \PendingCell \\
iPhos & \PendingCell & \PendingCell & \PendingCell & \PendingCell & \PendingCell \\
Ketone--isocyanide amide & \PendingCell & \PendingCell & \PendingCell & \PendingCell & \PendingCell \\
Ketone Ugi 4-CR & \PendingCell & \PendingCell & \PendingCell & \PendingCell & \PendingCell \\
Maleate & \PendingCell & \PendingCell & \PendingCell & \PendingCell & \PendingCell \\
O-esterification & \PendingCell & \PendingCell & \PendingCell & \PendingCell & \PendingCell \\
Passerini & \PendingCell & \PendingCell & \PendingCell & \PendingCell & \PendingCell \\
Preassembled thiol-yne & \PendingCell & \PendingCell & \PendingCell & \PendingCell & \PendingCell \\
Reductive amination & \PendingCell & \PendingCell & \PendingCell & \PendingCell & \PendingCell \\
Thiolactone & \PendingCell & \PendingCell & \PendingCell & \PendingCell & \PendingCell \\
\bottomrule\end{tabular}\end{table}
\begin{table}[H]
\caption{\textbf{Independent training-seed outcomes.} Each column denotes an independent training run, with identifiers \PendingValue{training seed IDs}. We report sample counts, paired differences and observed ranges for each method and inference arm. Sampling noise from one checkpoint does not constitute an independent training seed.}
\label{tab:22-seeds}
\centering\footnotesize\setlength{\tabcolsep}{3pt}\renewcommand{\arraystretch}{1.12}
\begin{tabular}{@{}>{\raggedright\arraybackslash}p{.30\textwidth}ccccc@{}}\toprule
Family & \shortstack{Seed A\\exact / $N$} & \shortstack{Seed B\\exact / $N$} & \shortstack{Seed C\\exact / $N$} & \shortstack{Mean $\pm$\\sample SD} & \shortstack{Paired control\\contrasts} \\\midrule
A3 amine--aldehyde--alkyne & \PendingCell & \PendingCell & \PendingCell & \PendingCell & \PendingCell \\
Acid--epoxide diester & \PendingCell & \PendingCell & \PendingCell & \PendingCell & \PendingCell \\
AEMA & \PendingCell & \PendingCell & \PendingCell & \PendingCell & \PendingCell \\
Aldehyde Ugi 3-CR & \PendingCell & \PendingCell & \PendingCell & \PendingCell & \PendingCell \\
Aldehyde Ugi 4-CR & \PendingCell & \PendingCell & \PendingCell & \PendingCell & \PendingCell \\
$\alpha$-Isocyanoester dihydroimidazole & \PendingCell & \PendingCell & \PendingCell & \PendingCell & \PendingCell \\
Amine alkylation & \PendingCell & \PendingCell & \PendingCell & \PendingCell & \PendingCell \\
Amine--epoxide & \PendingCell & \PendingCell & \PendingCell & \PendingCell & \PendingCell \\
Aryl reductive amination & \PendingCell & \PendingCell & \PendingCell & \PendingCell & \PendingCell \\
Aza-Michael acrylamide & \PendingCell & \PendingCell & \PendingCell & \PendingCell & \PendingCell \\
Aza-Michael acrylate & \PendingCell & \PendingCell & \PendingCell & \PendingCell & \PendingCell \\
Disulfide Michael & \PendingCell & \PendingCell & \PendingCell & \PendingCell & \PendingCell \\
Epoxide O-acylation & \PendingCell & \PendingCell & \PendingCell & \PendingCell & \PendingCell \\
iPhos & \PendingCell & \PendingCell & \PendingCell & \PendingCell & \PendingCell \\
Ketone--isocyanide amide & \PendingCell & \PendingCell & \PendingCell & \PendingCell & \PendingCell \\
Ketone Ugi 4-CR & \PendingCell & \PendingCell & \PendingCell & \PendingCell & \PendingCell \\
Maleate & \PendingCell & \PendingCell & \PendingCell & \PendingCell & \PendingCell \\
O-esterification & \PendingCell & \PendingCell & \PendingCell & \PendingCell & \PendingCell \\
Passerini & \PendingCell & \PendingCell & \PendingCell & \PendingCell & \PendingCell \\
Preassembled thiol-yne & \PendingCell & \PendingCell & \PendingCell & \PendingCell & \PendingCell \\
Reductive amination & \PendingCell & \PendingCell & \PendingCell & \PendingCell & \PendingCell \\
Thiolactone & \PendingCell & \PendingCell & \PendingCell & \PendingCell & \PendingCell \\
\bottomrule\end{tabular}\end{table}

\subsection{Matched controls, native baselines and the three-family comparison}
\begin{table}[H]
\caption{\textbf{Matched controls and three-family comparison.} Each comparison specifies compatible programs, splits, attempt budgets and verification. Unsupported native chemistry is unassessed,
not a zero score.}
\label{tab:22-controls}\centering\scriptsize
\begin{tabular}{@{}p{.43\textwidth}cccc@{}}\toprule
Arm / comparison & Exact L1 & Distinct novel & Quality axes & Cost\\\midrule
Conditioned shared model & \PendingCell & \PendingCell & \PendingCell & \PendingCell\\
Independently trained shared-null/post-hoc & \PendingCell & \PendingCell & \PendingCell & \PendingCell\\
Independently trained cyclic-label control & \PendingCell & \PendingCell & \PendingCell & \PendingCell\\
Finite TRAIN catalogue and learned selector & \PendingCell & \PendingCell & \PendingCell & \PendingCell\\
FACT, matched and generous capacity & \PendingCell & \PendingCell & \PendingCell & \PendingCell\\
DeFoG / GenMol-SAFE / RGFN, qualified scope & \PendingCell & \PendingCell & \PendingCell & \PendingCell\\
Three-family versus 22-family training, common programs & \PendingCell & \PendingCell & \PendingCell & \PendingCell\\
\bottomrule\end{tabular}\end{table}
Control comparisons use paired request distributions and report differences within each training
seed, observed ranges and uncertainty. Training exposure, checkpoint selection, decoder opportunity
and source/planner budgets define the matched comparison. We distinguish supported native chemistry
from unavailable baseline coverage. A wrong reaction label supplied alongside correct roles and
reaction cores tests the incremental information in that label; a retrained null model additionally
removes the corresponding training signal.

\subsection{Decoder and structural-repair attribution}
\begin{table}[H]
\caption{\textbf{Fixed-weight inference ablations.} Original, proposed and selected graphs share the same saved requests and model logits. Larger candidate-pool comparisons are separated from matched-cost controls; fixed-weight changes measure inference contributions.}
\label{tab:22-decoder}\centering\scriptsize
\begin{tabular}{@{}p{.43\textwidth}cccc@{}}\toprule
Inference arm & Exact L1 & Design checks & Diversity & Added cost\\\midrule
Raw readout & \PendingCell & \PendingCell & \PendingCell & \PendingCell\\
Constrained topology / core readout & \PendingCell & \PendingCell & \PendingCell & \PendingCell\\
Checked component completion & \PendingCell & \PendingCell & \PendingCell & \PendingCell\\
Scoped head-survival proposals & \PendingCell & \PendingCell & \PendingCell & \PendingCell\\
Joint role atom/bond motif proposals & \PendingCell & \PendingCell & \PendingCell & \PendingCell\\
Final ring-system topology proposals & \PendingCell & \PendingCell & \PendingCell & \PendingCell\\
Combined repairs, original candidates retained & \PendingCell & \PendingCell & \PendingCell & \PendingCell\\
Combined repairs + declared selection rule & \PendingCell & \PendingCell & \PendingCell & \PendingCell\\
Equal-cost control / no TRAIN-motif prior & \PendingCell & \PendingCell & \PendingCell & \PendingCell\\
\bottomrule\end{tabular}\end{table}
The inference analysis separates readout, component completion, local chemical proposals, ring-system
proposals and population selection. We retain original, proposed and selected graphs, with atom,
bond and topology edits, repeated-arm identities, rejected proposals and search limits. Every design
predicate is evaluated on the final graph because a later construction can undo an earlier correction.
Cost includes duplicate and failed proposals, source checks and selection. TRAIN local motifs used
as proposal priors are identified separately from independent evaluation references.

\subsection{Structural plausibility and design-condition agreement}
\begin{table}[H]
\caption{\textbf{Structural assessment by family.} All-request outcomes are accompanied by a separate matched source-positive control population. Each rule has applicable, assessed, pass, fail and abstain counts; head and ring agreement are distinct endpoints. Exact-plus-design yield requires every applicable check to be assessed and pass; abstentions and unknowns cannot enter its numerator, and all requests remain in its denominator; it is not a universal quality score. Context flags are not automatic chemical failures. Control disagreements and sensitivity limits define the scope of interpretation.}
\label{tab:22-structure}
\centering\footnotesize\setlength{\tabcolsep}{3pt}\renewcommand{\arraystretch}{1.12}
\begin{tabular}{@{}>{\raggedright\arraybackslash}p{.27\textwidth}cccccc@{}}
\toprule
Family & \shortstack{Exact + design\\pass / $N$} & \shortstack{Represent.\\failures} & \shortstack{Chemical\\alerts} & \shortstack{Context\\flags} & \shortstack{Head / ring\\agreement} & \shortstack{Abstain /\\unassessed} \\
\midrule
A3 amine--aldehyde--alkyne & \PendingCell & \PendingCell & \PendingCell & \PendingCell & \PendingCell & \PendingCell \\
Acid--epoxide diester & \PendingCell & \PendingCell & \PendingCell & \PendingCell & \PendingCell & \PendingCell \\
AEMA & \PendingCell & \PendingCell & \PendingCell & \PendingCell & \PendingCell & \PendingCell \\
Aldehyde Ugi 3-CR & \PendingCell & \PendingCell & \PendingCell & \PendingCell & \PendingCell & \PendingCell \\
Aldehyde Ugi 4-CR & \PendingCell & \PendingCell & \PendingCell & \PendingCell & \PendingCell & \PendingCell \\
$\alpha$-Isocyanoester dihydroimidazole & \PendingCell & \PendingCell & \PendingCell & \PendingCell & \PendingCell & \PendingCell \\
Amine alkylation & \PendingCell & \PendingCell & \PendingCell & \PendingCell & \PendingCell & \PendingCell \\
Amine--epoxide & \PendingCell & \PendingCell & \PendingCell & \PendingCell & \PendingCell & \PendingCell \\
Aryl reductive amination & \PendingCell & \PendingCell & \PendingCell & \PendingCell & \PendingCell & \PendingCell \\
Aza-Michael acrylamide & \PendingCell & \PendingCell & \PendingCell & \PendingCell & \PendingCell & \PendingCell \\
Aza-Michael acrylate & \PendingCell & \PendingCell & \PendingCell & \PendingCell & \PendingCell & \PendingCell \\
Disulfide Michael & \PendingCell & \PendingCell & \PendingCell & \PendingCell & \PendingCell & \PendingCell \\
Epoxide O-acylation & \PendingCell & \PendingCell & \PendingCell & \PendingCell & \PendingCell & \PendingCell \\
iPhos & \PendingCell & \PendingCell & \PendingCell & \PendingCell & \PendingCell & \PendingCell \\
Ketone--isocyanide amide & \PendingCell & \PendingCell & \PendingCell & \PendingCell & \PendingCell & \PendingCell \\
Ketone Ugi 4-CR & \PendingCell & \PendingCell & \PendingCell & \PendingCell & \PendingCell & \PendingCell \\
Maleate & \PendingCell & \PendingCell & \PendingCell & \PendingCell & \PendingCell & \PendingCell \\
O-esterification & \PendingCell & \PendingCell & \PendingCell & \PendingCell & \PendingCell & \PendingCell \\
Passerini & \PendingCell & \PendingCell & \PendingCell & \PendingCell & \PendingCell & \PendingCell \\
Preassembled thiol-yne & \PendingCell & \PendingCell & \PendingCell & \PendingCell & \PendingCell & \PendingCell \\
Reductive amination & \PendingCell & \PendingCell & \PendingCell & \PendingCell & \PendingCell & \PendingCell \\
Thiolactone & \PendingCell & \PendingCell & \PendingCell & \PendingCell & \PendingCell & \PendingCell \\
\midrule
Equal-family macro & \PendingCell & \PendingCell & \PendingCell & \PendingCell & \PendingCell & \PendingCell \\
Pooled all requests & \PendingCell & \PendingCell & \PendingCell & \PendingCell & \PendingCell & \PendingCell \\
Worst-family result & \PendingCell & \PendingCell & \PendingCell & \PendingCell & \PendingCell & \PendingCell \\
\bottomrule
\end{tabular}
\end{table}

Structural assessment records sanitization, connectedness, radical/charge/vocabulary support,
chemical alerts, retained head sites, carbon-domain sizes and branching, head-to-domain distances
and connected ring systems. Alert definitions are fixed before candidate review. No universal
tail-length, logP or pKa cutoff is imposed. Graph cycle rank and connected ring systems are distinct
from SSSR rings and serialization-dependent fundamental cycles; incompatible basis comparisons
abstain. Visual assessment uses a method-masked sample from every family and matched controls,
with sampling seed \PendingValue{visual sampling seed} and reviewers \PendingValue{reviewer type}.
Graph-specific reasons, uncertainty and depiction corrections accompany the review labels.

\begin{table}[H]
\caption{\textbf{Head, ring and TRAIN-support diagnostics by family.}
Head retention uses the declared family-scope request denominator, including charged/ambiguous
abstentions. Cycle allocation uses all requests, including matching zero closures; variable ring-size
agreement uses comparable cyclic requests, with unassessed counts. Atom coverage averages assessed
products. Exact-plus-observed requires product and all component features, divided by all requests.
TRAIN support is descriptive and never substitutes for independent quality evidence.}
\label{tab:22-conditions}\centering\footnotesize
\setlength{\tabcolsep}{3pt}\renewcommand{\arraystretch}{1.12}
\begin{tabular}{@{}>{\raggedright\arraybackslash}p{.30\textwidth}ccccc@{}}\toprule
Family & \shortstack{Head\\retention} & \shortstack{Role cycle\\allocation} & \shortstack{Ring-size\\agreement} & \shortstack{Atom-env.\\coverage} & \shortstack{Exact + observed\\product/comps. / $N$}\\\midrule
A3 amine--aldehyde--alkyne & \PendingCell & \PendingCell & \PendingCell & \PendingCell & \PendingCell \\
Acid--epoxide diester & \PendingCell & \PendingCell & \PendingCell & \PendingCell & \PendingCell \\
AEMA & \PendingCell & \PendingCell & \PendingCell & \PendingCell & \PendingCell \\
Aldehyde Ugi 3-CR & \PendingCell & \PendingCell & \PendingCell & \PendingCell & \PendingCell \\
Aldehyde Ugi 4-CR & \PendingCell & \PendingCell & \PendingCell & \PendingCell & \PendingCell \\
$\alpha$-Isocyanoester dihydroimidazole & \PendingCell & \PendingCell & \PendingCell & \PendingCell & \PendingCell \\
Amine alkylation & \PendingCell & \PendingCell & \PendingCell & \PendingCell & \PendingCell \\
Amine--epoxide & \PendingCell & \PendingCell & \PendingCell & \PendingCell & \PendingCell \\
Aryl reductive amination & \PendingCell & \PendingCell & \PendingCell & \PendingCell & \PendingCell \\
Aza-Michael acrylamide & \PendingCell & \PendingCell & \PendingCell & \PendingCell & \PendingCell \\
Aza-Michael acrylate & \PendingCell & \PendingCell & \PendingCell & \PendingCell & \PendingCell \\
Disulfide Michael & \PendingCell & \PendingCell & \PendingCell & \PendingCell & \PendingCell \\
Epoxide O-acylation & \PendingCell & \PendingCell & \PendingCell & \PendingCell & \PendingCell \\
iPhos & \PendingCell & \PendingCell & \PendingCell & \PendingCell & \PendingCell \\
Ketone--isocyanide amide & \PendingCell & \PendingCell & \PendingCell & \PendingCell & \PendingCell \\
Ketone Ugi 4-CR & \PendingCell & \PendingCell & \PendingCell & \PendingCell & \PendingCell \\
Maleate & \PendingCell & \PendingCell & \PendingCell & \PendingCell & \PendingCell \\
O-esterification & \PendingCell & \PendingCell & \PendingCell & \PendingCell & \PendingCell \\
Passerini & \PendingCell & \PendingCell & \PendingCell & \PendingCell & \PendingCell \\
Preassembled thiol-yne & \PendingCell & \PendingCell & \PendingCell & \PendingCell & \PendingCell \\
Reductive amination & \PendingCell & \PendingCell & \PendingCell & \PendingCell & \PendingCell \\
Thiolactone & \PendingCell & \PendingCell & \PendingCell & \PendingCell & \PendingCell \\
\bottomrule\end{tabular}\end{table}

\subsection{Lipid realism independent of decoder priors}
\begin{table}[H]
\caption{\textbf{Reference-distribution assessment.} We compare source-held-out empirical lipids and family-matched controls as separate reference populations. Descriptor distance is reported feature-wise, including tails, not collapsed to an arbitrary acceptability cutoff. Valid and unique counts, group splits, uncertainty and missing reference support accompany each comparison.}
\label{tab:22-realism}
\centering\footnotesize\setlength{\tabcolsep}{3pt}\renewcommand{\arraystretch}{1.12}
\begin{tabular}{@{}>{\raggedright\arraybackslash}p{.30\textwidth}ccccc@{}}
\toprule
Family & \shortstack{Fingerprint\\prec. / cov.} & \shortstack{Descriptor\\prec. / cov.} & \shortstack{Grouped\\C2ST AUC} & \shortstack{Descriptor\\distance} & \shortstack{Reference /\\generated $n$} \\
\midrule
A3 amine--aldehyde--alkyne & \PendingCell & \PendingCell & \PendingCell & \PendingCell & \PendingCell \\
Acid--epoxide diester & \PendingCell & \PendingCell & \PendingCell & \PendingCell & \PendingCell \\
AEMA & \PendingCell & \PendingCell & \PendingCell & \PendingCell & \PendingCell \\
Aldehyde Ugi 3-CR & \PendingCell & \PendingCell & \PendingCell & \PendingCell & \PendingCell \\
Aldehyde Ugi 4-CR & \PendingCell & \PendingCell & \PendingCell & \PendingCell & \PendingCell \\
$\alpha$-Isocyanoester dihydroimidazole & \PendingCell & \PendingCell & \PendingCell & \PendingCell & \PendingCell \\
Amine alkylation & \PendingCell & \PendingCell & \PendingCell & \PendingCell & \PendingCell \\
Amine--epoxide & \PendingCell & \PendingCell & \PendingCell & \PendingCell & \PendingCell \\
Aryl reductive amination & \PendingCell & \PendingCell & \PendingCell & \PendingCell & \PendingCell \\
Aza-Michael acrylamide & \PendingCell & \PendingCell & \PendingCell & \PendingCell & \PendingCell \\
Aza-Michael acrylate & \PendingCell & \PendingCell & \PendingCell & \PendingCell & \PendingCell \\
Disulfide Michael & \PendingCell & \PendingCell & \PendingCell & \PendingCell & \PendingCell \\
Epoxide O-acylation & \PendingCell & \PendingCell & \PendingCell & \PendingCell & \PendingCell \\
iPhos & \PendingCell & \PendingCell & \PendingCell & \PendingCell & \PendingCell \\
Ketone--isocyanide amide & \PendingCell & \PendingCell & \PendingCell & \PendingCell & \PendingCell \\
Ketone Ugi 4-CR & \PendingCell & \PendingCell & \PendingCell & \PendingCell & \PendingCell \\
Maleate & \PendingCell & \PendingCell & \PendingCell & \PendingCell & \PendingCell \\
O-esterification & \PendingCell & \PendingCell & \PendingCell & \PendingCell & \PendingCell \\
Passerini & \PendingCell & \PendingCell & \PendingCell & \PendingCell & \PendingCell \\
Preassembled thiol-yne & \PendingCell & \PendingCell & \PendingCell & \PendingCell & \PendingCell \\
Reductive amination & \PendingCell & \PendingCell & \PendingCell & \PendingCell & \PendingCell \\
Thiolactone & \PendingCell & \PendingCell & \PendingCell & \PendingCell & \PendingCell \\
\midrule
Equal-family macro & \PendingCell & \PendingCell & \PendingCell & \PendingCell & \PendingCell \\
Pooled all requests & \PendingCell & \PendingCell & \PendingCell & \PendingCell & \PendingCell \\
Worst-family result & \PendingCell & \PendingCell & \PendingCell & \PendingCell & \PendingCell \\
\bottomrule
\end{tabular}
\end{table}

\noindent\textbf{Aggregation.} AUC and distribution distances are not pooled success ratios. We
report equal-family macro means, recomputed pooled estimates and worst-family values separately,
with the direction of each metric stated. Unsupported reference populations remain unassessed;
missing-family counts accompany each aggregate.

Descriptor comparisons cover molecular size, lipophilicity, polar surface area, hydrogen-bond
counts, flexibility, rings and aromaticity, charge, heteroatom composition, branching, unsaturation
and graph-based lipid regions. We report feature-wise reference distributions and normalized
Wasserstein distances. Fingerprint precision and coverage, source-grouped classifier AUC and TRAIN
feature membership quantify different properties. Source-grouped splits exclude duplicate identities
across classifier partitions and do not supply source identifiers as features. No descriptor or
similarity metric establishes biological utility.

\subsection{Exact definitions of the quality and realism metrics}
\label{app:22-metrics}
The metrics below distinguish structural conditions, reference membership and distributional
similarity. Reference populations, constitutional identity, implementations and random seeds are fixed
for each comparison. A metric's definition does not imply that a generated population satisfies it.

\begingroup\footnotesize\setlength{\tabcolsep}{4pt}\renewcommand{\arraystretch}{1.12}
\begin{longtable}{@{}>{\raggedright\arraybackslash}p{.30\textwidth}>{\raggedright\arraybackslash}p{.64\textwidth}@{}}
\caption{\textbf{Quality and realism metric definitions.} Each metric has an explicit numerator, denominator and interpretation. Structural checks and reference-support
metrics are not interchangeable.}\label{tab:22-metric-definitions}\\
\toprule Metric & Exact computation / limitation\\\midrule\endfirsthead
\multicolumn{2}{l}{Table~\ref{tab:22-metric-definitions}, continued}\\\toprule
Metric & Exact computation / limitation\\\midrule\endhead
\bottomrule\endfoot
Valid molecules / all requests (\%) & RDKit parsing/sanitization; failure remains in all-request denominator. \\[1.5pt]
Connected molecules / all requests (\%) & Exactly one connected molecular graph; distinct from validity. \\[1.5pt]
Exact L1 + all qualified design checks / requests (\%) & Source-exact final graph and every applicable qualified axis assessed/pass. Unknowns and abstentions do not pass. \\[1.5pt]
Molecules with a justified chemical alert / requests (\%) & Frozen narrow chemical-alert motifs, with per-rule counts and unassessed requests reported separately. This is not a complete stability screen. \\[1.5pt]
Molecules requiring chemistry-context review / requests (\%) & Frozen context flags are descriptive, not automatic failures. \\[1.5pt]
Required head retained / family-scope requests (\%) & Pass / prespecified family-scope requests, including charged or ambiguous-origin abstentions. Only neutral aldehyde-Ugi4 has a qualified predicate; other families remain unassessed. \\[1.5pt]
Correct role-local cycle allocation / requests (\%) & Match declared variable closures per role, including zero; cross-origin edges must also match. All requests remain in the denominator; unassessed cases are reported separately. \\[1.5pt]
Requested ring-size agreement / comparable requests (\%) & Each observed fundamental-cycle length must belong to its role's allowed-size set, under a comparable tree basis. Not a size-multiset, SSSR or 3D-geometry match. Inapplicable and unassessed counts are reported separately. \\[1.5pt]
TRAIN atom-environment coverage (\%) & Fraction of atoms with reference-observed local environments, averaged over assessed products. Component-role summaries are reported separately. \\[1.5pt]
Ring-system TRAIN support / assessed products (\%) & Among valid connected products with a family TRAIN reference, fraction with no unknown connected ring-bond signature. Ring-free products pass vacuously; their counts and missing/invalid cases are reported separately. Component roles are separate. \\[1.5pt]
Exact L1 + product/component features observed / requests (\%) & Exact, unambiguously decomposed product and every precursor must have all atom environments and connected ring-bond signatures observed in their family product or family/source-role TRAIN reference. The denominator is all requests; absent references never pass. \\[1.5pt]
Fingerprint precision (\%) & Fraction of unique generated molecules inside the union of reference k-nearest-neighbor balls in ECFP4/Tanimoto distance. \\[1.5pt]
Fingerprint coverage (\%) & Fraction of reference-centered neighborhoods hit by at least one unique generated molecule. \\[1.5pt]
Fingerprint recall (\%) & Fraction of references inside generated-centered neighborhoods. Distinct from coverage; small generated sets can inflate neighborhood radii. \\[1.5pt]
Distinct reference-supported products / requests (\%) & Unique generated identities inside reference neighborhoods divided by all requests; exposes invalidity and duplication. \\[1.5pt]
Nearest-reference Tanimoto similarity & Mean nearest-reference similarity over unique generated identities; descriptive, not a novelty or quality objective. \\[1.5pt]
Internal ECFP4 diversity & Mean Tanimoto distance between distinct generated identities, interpreted alongside validity, realism and component diversity. \\[1.5pt]
Descriptor distribution distance, per feature & One-dimensional Wasserstein distance / reference IQR; near-zero fallback: max(reference population SD,1). Connected-request observations retain duplicates. Smaller means closer to reference, not better chemistry. \\[1.5pt]
Descriptor-space precision and coverage & Precision measures generated points in reference-centered descriptor neighborhoods; coverage measures reference neighborhoods reached by generated points. Scaling and estimator: \PendingValue{descriptor neighborhood estimator}; source-held-out reference: \PendingValue{reference cohort}. \\[1.5pt]
Source-grouped two-sample classifier AUC & AUC distinguishes generated from reference molecules across held-out source groups. Classifier and sampling: \PendingValue{C2ST estimator}; grouping and duplicate controls: \PendingValue{C2ST split}. \\[1.5pt]
\end{longtable}\endgroup

\begin{table}[H]
\caption{\textbf{Lipid descriptor distributions.}
The 24 descriptors are compared feature by feature on connected-request observations, retaining
duplicates. We report means, SDs and distribution tails separately from model-seed uncertainty,
alongside normalized Wasserstein distance. $W_1/s_R$ denotes reference-scaled one-dimensional Wasserstein distance.
Each family/reference cohort has a separate panel; no arbitrary acceptance threshold is imposed.}
\label{tab:22-descriptors}\centering\footnotesize
\setlength{\tabcolsep}{3pt}\renewcommand{\arraystretch}{1.15}
\begin{tabular}{@{}>{\raggedright\arraybackslash}p{.35\textwidth}>{\raggedright\arraybackslash}p{.16\textwidth}cccc@{}}\toprule
Descriptor & Unit & Reference & Raw & Final & $W_1/s_R$\\\midrule
Heavy atoms & count & \PendingCell & \PendingCell & \PendingCell & \PendingCell \\
Molecular weight & Da & \PendingCell & \PendingCell & \PendingCell & \PendingCell \\
Calculated logP & unitless & \PendingCell & \PendingCell & \PendingCell & \PendingCell \\
Topological polar surface area & \AA$^2$ & \PendingCell & \PendingCell & \PendingCell & \PendingCell \\
Hydrogen-bond donors & count & \PendingCell & \PendingCell & \PendingCell & \PendingCell \\
Hydrogen-bond acceptors & count & \PendingCell & \PendingCell & \PendingCell & \PendingCell \\
Rotatable bonds & count & \PendingCell & \PendingCell & \PendingCell & \PendingCell \\
Rings & count & \PendingCell & \PendingCell & \PendingCell & \PendingCell \\
Aromatic rings & count & \PendingCell & \PendingCell & \PendingCell & \PendingCell \\
Fraction sp$^3$ carbons & fraction & \PendingCell & \PendingCell & \PendingCell & \PendingCell \\
Net formal charge & elementary charge & \PendingCell & \PendingCell & \PendingCell & \PendingCell \\
Sum of absolute formal charges & elementary charge & \PendingCell & \PendingCell & \PendingCell & \PendingCell \\
Carbon fraction & fraction & \PendingCell & \PendingCell & \PendingCell & \PendingCell \\
Nitrogen atoms & count & \PendingCell & \PendingCell & \PendingCell & \PendingCell \\
Oxygen atoms & count & \PendingCell & \PendingCell & \PendingCell & \PendingCell \\
Sulfur atoms & count & \PendingCell & \PendingCell & \PendingCell & \PendingCell \\
Phosphorus atoms & count & \PendingCell & \PendingCell & \PendingCell & \PendingCell \\
Branching heavy atoms & count & \PendingCell & \PendingCell & \PendingCell & \PendingCell \\
Nonaromatic unsaturated bonds & count & \PendingCell & \PendingCell & \PendingCell & \PendingCell \\
Maximum polar-root distance & graph bonds & \PendingCell & \PendingCell & \PendingCell & \PendingCell \\
Heuristic head fraction & fraction & \PendingCell & \PendingCell & \PendingCell & \PendingCell \\
Heuristic interface fraction & fraction & \PendingCell & \PendingCell & \PendingCell & \PendingCell \\
Heuristic tail fraction & fraction & \PendingCell & \PendingCell & \PendingCell & \PendingCell \\
Heuristic tail branching atoms & count & \PendingCell & \PendingCell & \PendingCell & \PendingCell \\
\bottomrule\end{tabular}
\par\smallskip\raggedright\scriptsize
Head/interface/tail fractions and polar-root distance use the descriptor implementation's graph
heuristic; they are not precursor-origin annotations or an independently qualified headgroup test.
Branching heavy atoms use heavy-atom degree at least three; carbon-domain branching in the
structural audit is a separate measurement. Formal charge is not ionization at physiological pH.
\end{table}

\subsection{Product and precursor diversity, novelty and reachability}
\begin{table}[H]
\caption{\textbf{Product and component diversity and novelty.} $N$ is all requests. $D_1$ and $D_2$ are Shannon and inverse-Simpson effective counts, reported per registered role. Novel-component endpoints require an unambiguous exact decomposition and include separate observation and distinct-product yields. TRAIN and full-declared-reference novelty are separate endpoints.}
\label{tab:22-diversity}
\centering\footnotesize\setlength{\tabcolsep}{3pt}\renewcommand{\arraystretch}{1.12}
\begin{tabular}{@{}>{\raggedright\arraybackslash}p{.30\textwidth}ccccc@{}}
\toprule
Family & \shortstack{Distinct\\exact / $N$} & \shortstack{Novel\\product / $N$} & \shortstack{Novel comp.\\product / $N$} & \shortstack{Component\\$D_1$ / $D_2$} & \shortstack{Novel distinct\\components} \\
\midrule
A3 amine--aldehyde--alkyne & \PendingCell & \PendingCell & \PendingCell & \PendingCell & \PendingCell \\
Acid--epoxide diester & \PendingCell & \PendingCell & \PendingCell & \PendingCell & \PendingCell \\
AEMA & \PendingCell & \PendingCell & \PendingCell & \PendingCell & \PendingCell \\
Aldehyde Ugi 3-CR & \PendingCell & \PendingCell & \PendingCell & \PendingCell & \PendingCell \\
Aldehyde Ugi 4-CR & \PendingCell & \PendingCell & \PendingCell & \PendingCell & \PendingCell \\
$\alpha$-Isocyanoester dihydroimidazole & \PendingCell & \PendingCell & \PendingCell & \PendingCell & \PendingCell \\
Amine alkylation & \PendingCell & \PendingCell & \PendingCell & \PendingCell & \PendingCell \\
Amine--epoxide & \PendingCell & \PendingCell & \PendingCell & \PendingCell & \PendingCell \\
Aryl reductive amination & \PendingCell & \PendingCell & \PendingCell & \PendingCell & \PendingCell \\
Aza-Michael acrylamide & \PendingCell & \PendingCell & \PendingCell & \PendingCell & \PendingCell \\
Aza-Michael acrylate & \PendingCell & \PendingCell & \PendingCell & \PendingCell & \PendingCell \\
Disulfide Michael & \PendingCell & \PendingCell & \PendingCell & \PendingCell & \PendingCell \\
Epoxide O-acylation & \PendingCell & \PendingCell & \PendingCell & \PendingCell & \PendingCell \\
iPhos & \PendingCell & \PendingCell & \PendingCell & \PendingCell & \PendingCell \\
Ketone--isocyanide amide & \PendingCell & \PendingCell & \PendingCell & \PendingCell & \PendingCell \\
Ketone Ugi 4-CR & \PendingCell & \PendingCell & \PendingCell & \PendingCell & \PendingCell \\
Maleate & \PendingCell & \PendingCell & \PendingCell & \PendingCell & \PendingCell \\
O-esterification & \PendingCell & \PendingCell & \PendingCell & \PendingCell & \PendingCell \\
Passerini & \PendingCell & \PendingCell & \PendingCell & \PendingCell & \PendingCell \\
Preassembled thiol-yne & \PendingCell & \PendingCell & \PendingCell & \PendingCell & \PendingCell \\
Reductive amination & \PendingCell & \PendingCell & \PendingCell & \PendingCell & \PendingCell \\
Thiolactone & \PendingCell & \PendingCell & \PendingCell & \PendingCell & \PendingCell \\
\midrule
Equal-family macro & \PendingCell & \PendingCell & \PendingCell & \PendingCell & \PendingCell \\
Pooled all requests & \PendingCell & \PendingCell & \PendingCell & \PendingCell & \PendingCell \\
Worst-family result & \PendingCell & \PendingCell & \PendingCell & \PendingCell & \PendingCell \\
\bottomrule
\end{tabular}
\end{table}

Product and per-role component distributions are summarized by unique counts, frequency tables,
largest-identity concentration and effective diversity. Novelty is evaluated separately against TRAIN
and the full declared reference. Absolute noninferiority comparisons preserve the baseline observation
denominators as well as unique and effective counts, including novel observations and distinct novel
identities. Global unique counts deduplicate across families. Catalogue inputs, admissible precursor
tuples and reachable unique products are separate quantities: a novel accepted precursor does not
by itself show that the product is unreachable through every catalogue decomposition.

\subsection{Generalization beyond development requests}
\begin{table}[H]
\caption{\textbf{Generalization and held-component recovery.} Each split has its own request denominator and independent training-seed outcomes. Unsupported cases remain unassessed. Size, closure and depth strata are reported separately. Zero-shot held-reaction-family behavior is a secondary stress test, not a hard gate.}
\label{tab:22-generalization}
\centering\footnotesize\setlength{\tabcolsep}{3pt}\renewcommand{\arraystretch}{1.12}
\begin{tabular}{@{}>{\raggedright\arraybackslash}p{.30\textwidth}ccccc@{}}
\toprule
Family & \shortstack{Held comp.\\yield / $N$} & \shortstack{Held identity\\coverage} & \shortstack{Component-\\disjoint L1} & \shortstack{Source-\\disjoint L1} & \shortstack{Size/repeat\\strata} \\
\midrule
A3 amine--aldehyde--alkyne & \PendingCell & \PendingCell & \PendingCell & \PendingCell & \PendingCell \\
Acid--epoxide diester & \PendingCell & \PendingCell & \PendingCell & \PendingCell & \PendingCell \\
AEMA & \PendingCell & \PendingCell & \PendingCell & \PendingCell & \PendingCell \\
Aldehyde Ugi 3-CR & \PendingCell & \PendingCell & \PendingCell & \PendingCell & \PendingCell \\
Aldehyde Ugi 4-CR & \PendingCell & \PendingCell & \PendingCell & \PendingCell & \PendingCell \\
$\alpha$-Isocyanoester dihydroimidazole & \PendingCell & \PendingCell & \PendingCell & \PendingCell & \PendingCell \\
Amine alkylation & \PendingCell & \PendingCell & \PendingCell & \PendingCell & \PendingCell \\
Amine--epoxide & \PendingCell & \PendingCell & \PendingCell & \PendingCell & \PendingCell \\
Aryl reductive amination & \PendingCell & \PendingCell & \PendingCell & \PendingCell & \PendingCell \\
Aza-Michael acrylamide & \PendingCell & \PendingCell & \PendingCell & \PendingCell & \PendingCell \\
Aza-Michael acrylate & \PendingCell & \PendingCell & \PendingCell & \PendingCell & \PendingCell \\
Disulfide Michael & \PendingCell & \PendingCell & \PendingCell & \PendingCell & \PendingCell \\
Epoxide O-acylation & \PendingCell & \PendingCell & \PendingCell & \PendingCell & \PendingCell \\
iPhos & \PendingCell & \PendingCell & \PendingCell & \PendingCell & \PendingCell \\
Ketone--isocyanide amide & \PendingCell & \PendingCell & \PendingCell & \PendingCell & \PendingCell \\
Ketone Ugi 4-CR & \PendingCell & \PendingCell & \PendingCell & \PendingCell & \PendingCell \\
Maleate & \PendingCell & \PendingCell & \PendingCell & \PendingCell & \PendingCell \\
O-esterification & \PendingCell & \PendingCell & \PendingCell & \PendingCell & \PendingCell \\
Passerini & \PendingCell & \PendingCell & \PendingCell & \PendingCell & \PendingCell \\
Preassembled thiol-yne & \PendingCell & \PendingCell & \PendingCell & \PendingCell & \PendingCell \\
Reductive amination & \PendingCell & \PendingCell & \PendingCell & \PendingCell & \PendingCell \\
Thiolactone & \PendingCell & \PendingCell & \PendingCell & \PendingCell & \PendingCell \\
\midrule
Equal-family macro & \PendingCell & \PendingCell & \PendingCell & \PendingCell & \PendingCell \\
Pooled all requests & \PendingCell & \PendingCell & \PendingCell & \PendingCell & \PendingCell \\
Worst-family result & \PendingCell & \PendingCell & \PendingCell & \PendingCell & \PendingCell \\
\bottomrule
\end{tabular}
\end{table}

Model, decoder and selection choices are fixed on development data before evaluation on the
declared held-out splits. We distinguish independent training seeds from sampling seeds and report
both request/layout identity overlap and graph-geometry overlap. Fresh noise on TRAIN-derived
layouts measures sampling variability, not held-out generalization. Uncertainty uses
\PendingValue{population-aware uncertainty estimator} to account for dependence introduced by
joint population selection.

\subsection{Upstream synthesis and terminal-evidence completeness}
\begin{table}[H]
\caption{\textbf{L1/L2/L3 evidence and dossier dispositions.} Products, unique family-role component constitutions, reaction steps and terminal leaves are distinct denominators. Complete dossiers require all admitted steps and leaves for the actual generated identity. L2/L3 unknowns, source conflicts, search limits and access failures are counted separately.}
\label{tab:22-routes}
\centering\footnotesize\setlength{\tabcolsep}{3pt}\renewcommand{\arraystretch}{1.12}
\begin{tabular}{@{}>{\raggedright\arraybackslash}p{.30\textwidth}ccccc@{}}
\toprule
Family & \shortstack{Exact L1\\products} & \shortstack{L2 closed /\\branches} & \shortstack{L3 qualified /\\leaves} & \shortstack{Complete\\dossiers / $N$} & \shortstack{Unresolved /\\censored} \\
\midrule
A3 amine--aldehyde--alkyne & \PendingCell & \PendingCell & \PendingCell & \PendingCell & \PendingCell \\
Acid--epoxide diester & \PendingCell & \PendingCell & \PendingCell & \PendingCell & \PendingCell \\
AEMA & \PendingCell & \PendingCell & \PendingCell & \PendingCell & \PendingCell \\
Aldehyde Ugi 3-CR & \PendingCell & \PendingCell & \PendingCell & \PendingCell & \PendingCell \\
Aldehyde Ugi 4-CR & \PendingCell & \PendingCell & \PendingCell & \PendingCell & \PendingCell \\
$\alpha$-Isocyanoester dihydroimidazole & \PendingCell & \PendingCell & \PendingCell & \PendingCell & \PendingCell \\
Amine alkylation & \PendingCell & \PendingCell & \PendingCell & \PendingCell & \PendingCell \\
Amine--epoxide & \PendingCell & \PendingCell & \PendingCell & \PendingCell & \PendingCell \\
Aryl reductive amination & \PendingCell & \PendingCell & \PendingCell & \PendingCell & \PendingCell \\
Aza-Michael acrylamide & \PendingCell & \PendingCell & \PendingCell & \PendingCell & \PendingCell \\
Aza-Michael acrylate & \PendingCell & \PendingCell & \PendingCell & \PendingCell & \PendingCell \\
Disulfide Michael & \PendingCell & \PendingCell & \PendingCell & \PendingCell & \PendingCell \\
Epoxide O-acylation & \PendingCell & \PendingCell & \PendingCell & \PendingCell & \PendingCell \\
iPhos & \PendingCell & \PendingCell & \PendingCell & \PendingCell & \PendingCell \\
Ketone--isocyanide amide & \PendingCell & \PendingCell & \PendingCell & \PendingCell & \PendingCell \\
Ketone Ugi 4-CR & \PendingCell & \PendingCell & \PendingCell & \PendingCell & \PendingCell \\
Maleate & \PendingCell & \PendingCell & \PendingCell & \PendingCell & \PendingCell \\
O-esterification & \PendingCell & \PendingCell & \PendingCell & \PendingCell & \PendingCell \\
Passerini & \PendingCell & \PendingCell & \PendingCell & \PendingCell & \PendingCell \\
Preassembled thiol-yne & \PendingCell & \PendingCell & \PendingCell & \PendingCell & \PendingCell \\
Reductive amination & \PendingCell & \PendingCell & \PendingCell & \PendingCell & \PendingCell \\
Thiolactone & \PendingCell & \PendingCell & \PendingCell & \PendingCell & \PendingCell \\
\midrule
Equal-family macro & \PendingCell & \PendingCell & \PendingCell & \PendingCell & \PendingCell \\
Pooled all requests & \PendingCell & \PendingCell & \PendingCell & \PendingCell & \PendingCell \\
Worst-family result & \PendingCell & \PendingCell & \PendingCell & \PendingCell & \PendingCell \\
\bottomrule
\end{tabular}
\end{table}

Route assessment follows the actual generated precursor identity. Method-masked adjudication of
the component union records source identifiers and locators, source/file hashes, reactive-site and
forward checks, terminal identity and availability dates, route depths, planner calls and elapsed time.
We report E0--E3 dispositions and distinguish raw from selection-conditioned coverage. An E2
assignment requires closure of the routes supporting its generated precursors. Evidence depth is
family-specific; detailed Ugi evidence is not transferred to another reaction family. L3 denotes
documentary terminal-material qualification rather than prospective procurement or synthesis.

\subsection{Architecture, theoretical assumptions and synthesis guidance}
\begin{table}[H]\caption{\textbf{Architecture comparisons.}
Comparisons use matched retraining and report training exposure, capacity and inference cost separately.}
\label{tab:22-architecture}\centering\small
\begin{tabular}{@{}p{.70\textwidth}l@{}}\toprule
Contrast / readouts & Result\\\midrule
Whole graph versus role-factorized, matched/generous capacity & \PendingCell\\
Input-only versus layerwise program conditioning & \PendingCell\\
Role/core, repeat, topology and chemistry-loss contributions & \PendingCell\\
Routed adapters/heads and gradient-conflict handling & \PendingCell\\
Per-family exact yield, held losses, diversity and dependence & \PendingCell\\
\bottomrule\end{tabular}\end{table}
\begin{table}[H]\caption{\textbf{Correspondence between theory and implementation.}}
\label{tab:22-theory}\centering\small
\begin{tabular}{@{}p{.70\textwidth}l@{}}\toprule
Obligation & Evidence\\\midrule
Bounded vocabularies/size/closures and actual production masks & \PendingCell\\
Reference objective versus implemented training reductions/losses & \PendingCell\\
Numerical sampler, endpoint readout, repair and selection maps & \PendingCell\\
Preserved invariants and properties requiring empirical evidence & \PendingCell\\
Independent-draw assumptions versus batch-coupled selection & \PendingCell\\
\bottomrule\end{tabular}\end{table}
\begin{table}[H]\caption{\textbf{Synthesis-guidance comparisons.}
We distinguish synthesis-guidance effects from a separately scoped potency diagnostic; all comparisons retain eligibility, abstention and budget accounting.}
\label{tab:22-guidance}\centering\small
\begin{tabular}{@{}p{.70\textwidth}l@{}}\toprule
Validation / contrast & Result\\\midrule
Nontrivial generated L1/L2/L3 dossiers and zero-guidance identity & \PendingCell\\
Matched unguided, post-hoc and in-trajectory synthesis contrast & \PendingCell\\
Eligibility, abstention, route yield, diversity and total budget & \PendingCell\\
Potency-only diagnostic with separate applicability policy & \PendingCell\\
\bottomrule\end{tabular}\end{table}
Architecture comparisons distinguish matched retraining from fixed-weight inference changes.
We relate each theoretical assumption to the implemented vocabulary, graph support, objective,
finite-step sampler and selection map. The ideal flow theorem does not certify a learned numerical
sampler, a TRAIN-derived repair or a batch-coupled selection procedure. Synthesis-guidance comparisons
separate unguided, post-hoc and in-trajectory arms at matched total budgets and require both a
zero-guidance identity check and auditable L1/L2/L3 dossiers. Outcomes outside the declared evidence
or applicability scope remain unassessed.

\subsection{Assembly, structural and synthesis-evidence outcomes}
\begin{figure}[H]\centering
\ResultPanel{a. Assembly and conditions}{Per-family raw and selected exact yield, request counts and training-seed variability.}\hfill
\ResultPanel{b. Structural assessment}{Chemical alerts, qualified condition failures and explicit abstentions.}\\[6pt]
\ResultPanel{c. Diversity and realism}{Product and role diversity, novelty and reference precision/coverage.}\hfill
\ResultPanel{d. Synthesis evidence}{L1, L2 and L3 dispositions and complete-dossier yield per request.}
\caption{\textbf{Assembly, molecular properties and synthesis evidence across reaction families.}
Assembly consistency, structural assessment, distributional similarity and evidence depth are
separate outcomes. Family-wise estimates retain their own denominators, comparator budgets and
training-seed uncertainty. Raw and selected populations are distinguished throughout.}
\label{fig:22-overview}
\end{figure}

\subsection{Representative structures and unresolved failures}
\begin{figure}[H]\centering
\ResultPanel{a. Seeded all-family atlas}{One or more prespecified random products per family, with exact precursors and evidence depth.}\hfill
\ResultPanel{b. Paired failure analysis}{Original and repaired graphs; edits, remaining concerns and unchanged source checks.}
\caption{\textbf{Structures and failure analysis across reaction families.} The atlas uses a fixed selection seed across all families. Atom origins, component novelty and L1/L2/L3 evidence are annotated separately. Random samples and targeted failure examples are distinguished. Molecular drawings represent exact graphs and do not establish three-dimensional stability.}
\label{fig:22-atlas}\end{figure}
The structure analysis pairs original and repaired graphs and includes oversized rings, malformed
cages, missing retained heads and joint atom/bond errors, alongside source-positive controls.
Graph and drawing hashes distinguish molecular changes from depiction corrections. Random samples
and examples selected to illustrate a specific failure are identified separately.

\subsection{Failure accounting and reproducibility}
The analysis retains failed requests and runs, unsupported chemistry, unavailable references,
source conflicts, null contrasts and search censoring. Source, configuration, checkpoint, policy and
split hashes link each result to its executable analysis and verification receipt. We distinguish
TRAIN-derived proposal and selection diagnostics from evaluation against independent references.
Passing limited structural checks does not establish complete chemical quality, stability, ionization
or delivery activity; context flags and unresolved cases remain part of the reported outcome.


\clearpage
\section{Theoretical analysis and proofs}
\label{app:theory}

The results establish separate properties of the construction. Bounded completeness concerns the
state representation. Catalogue reachability concerns exact precursor inputs. The population
argument connects a reference denoising loss to an ideal continuous-time law. Budget and route
results concern computational procedures. None establishes experimental synthesis success or
positive learned probability for every representable graph.

\subsection{Graph states and program compatibility}
\label{app:state-support}

The graph space $\mathfrak G_{N,K_{\rm cl}}$ is defined in
Section~\ref{sec:mathematical-objects}, with finite atom and bond alphabets and integers $N\geq1$,
$K_{\rm cl}\geq0$. Molecular validity is an additional predicate on that space.
The raw space $\overline{\mathfrak X}_{N,K_{\rm cl}}$ contains padded categorical states.

\begin{definition}[Well-formed sparse state]
\label{def:sparse-state}
A state belongs to $\mathfrak X_{N,K_{\rm cl}}$ when, for some $1\leq n\leq N$, its active atoms
are exactly positions $1,\ldots,n$ with labels in $\Sigma_{\rm atom}$; position 1 is the sole root
with null parent and parent-bond states; and every $2\leq i\leq n$ has an earlier parent
$1\leq\pi_i<i$ and a bond label in $\Sigma_{\rm bond}$. Each active closure has endpoints
$1\leq e^L<e^R\leq n$, a bond label in $\Sigma_{\rm bond}$, and an edge distinct from every tree
edge and other closure. Active closures precede unused slots and are sorted by endpoint pair.
All unused atom, pointer and bond coordinates carry their designated null values.
\end{definition}
The decoder reconstructs the labeled edges of a well-formed state and returns $\bot_{\rm dec}$
otherwise. It is therefore a total map into $\mathfrak G_{N,K_{\rm cl}}\cup\{\bot_{\rm dec}\}$.

\begin{lemma}[Tree residual count]
\label{lem:tree-closure}
For a finite nonempty connected graph $G=(V,E)$ and any spanning tree $T$,
$|E\setminus E(T)|=|E|-|V|+1=\beta(G)$.
\end{lemma}
\begin{proof}
The tree contains $|V|-1$ edges, all in $E$. Subtracting this count from $|E|$ gives the result,
including the single-vertex case.
\end{proof}

\begin{theorem}[Bounded completeness of the sparse representation]
\label{thm:bounded-representation}
For the spaces in Definition~\ref{def:sparse-state}, the restriction
$\operatorname{Dec}:\mathfrak X_{N,K_{\rm cl}}\to\mathfrak G_{N,K_{\rm cl}}$ is well defined
and surjective. Thus every well-formed state decodes to a graph in the declared support, and every
graph in that support has an encoding that decodes to its constitutional isomorphism class.
\end{theorem}
\begin{proof}
We first reconstruct a graph from a state, then construct a state for an arbitrary graph.
In a well-formed state with $n$ active atoms, every parent pointer strictly decreases its index.
Following parents must reach position 1, so the $n-1$ parent edges form a connected tree.
Adding $k\leq K_{\rm cl}$ distinct, non-loop closure edges preserves connectedness and gives cycle
rank $(n-1+k)-n+1=k$. The label and padding conditions ensure the decoded graph belongs to
$\mathfrak G_{N,K_{\rm cl}}$.

Conversely, choose a representative of any $G\in\mathfrak G_{N,K_{\rm cl}}$ and a rooted spanning
tree. Enumerate vertices in a traversal that places every parent before its children. Store their
labels, parent indices and tree-bond labels. Lemma~\ref{lem:tree-closure} leaves exactly
$\beta(G)\leq K_{\rm cl}$ residual edges; store each with increasing endpoint indices and sort
these triples. Pad all unused coordinates with null states. This state satisfies
Definition~\ref{def:sparse-state}, and the vertex enumeration is a label-preserving isomorphism
between its decoded graph and $G$.
\end{proof}

This result concerns the unrestricted representation. For program-specific alphabets
$\mathcal X_i(P)$ and fixed values $x_{P,i}$, define
\begin{equation}
 \Omega_P=\{x\in\overline{\mathfrak X}_{N,K_{\rm cl}}:
 x_i\in\mathcal X_i(P)\ \forall i,\quad x_i=x_{P,i}\ \forall i\in\mathcal F(P)\}.
 \label{eq:program-state-set}
\end{equation}
Its nonempty well-formed intersection defines
$\mathcal P_{\rm adm}=\{P:\mathfrak X_{N,K_{\rm cl}}\cap\Omega_P\neq\varnothing\}$.
Nonemptiness alone establishes neither chemical validity nor satisfaction of every morphology
constraint. Corpus admission additionally checks the annotated encoding, program alignment and
exact decode round trip. Representational completeness makes no claim that every graph is
encodable under every $P$.

\subsection{Exact assembly and catalogue reachability}
\label{app:catalogue-reachability}

Definition~\ref{def:exact-assembly} distinguishes the complete input preimage
$\mathcal E_c^\star(G)$ from recovered input--trace pairs $\widehat{\mathcal E}_c(G)$.
Fix a registered chemistry $c$ and role-specific catalogues $\mathcal C=(\mathcal C_r)_r$. Let
$\mathfrak C_c(\mathcal C)\subseteq\mathfrak C_c$ contain the admissible precursor inputs whose
every role-assigned component belongs to its catalogue. The reachable product set is
\begin{equation}
 \mathcal G_{\rm cat}(c,\mathcal C)
 =\bigcup_{\mathbf C\in\mathfrak C_c(\mathcal C)}\mathcal F_c(\mathbf C).
 \label{eq:catalogue-reachable}
\end{equation}
This definition includes all registered outcomes of the admissible inputs, not only products found
by a bounded search.

\begin{proposition}[Finite-component reachability]
\label{prop:catalogue-bound}
If $\mathfrak C_c(\mathcal C)$ and each $\mathcal F_c(\mathbf C)$ are finite, then
\begin{equation}
 |\mathcal G_{\rm cat}(c,\mathcal C)|
 \leq\sum_{\mathbf C\in\mathfrak C_c(\mathcal C)}|\mathcal F_c(\mathbf C)|
 \leq m_{\max,c}(\mathcal C)\,|\mathfrak C_c(\mathcal C)|,
\end{equation}
where $m_{\max,c}(\mathcal C)$ is the largest outcome count, or zero for an empty input set.
For every $G\in\mathfrak G_{N,K_{\rm cl}}$,
\begin{equation}
 G\in\mathcal G_{\rm cat}(c,\mathcal C)
 \quad\Longleftrightarrow\quad
 \mathcal E_c^\star(G)\cap\mathfrak C_c(\mathcal C)\neq\varnothing.
\end{equation}
\end{proposition}
\begin{proof}
The cardinality of a finite union is at most the sum of its member-set cardinalities; bounding each
summand by its maximum proves both inequalities. If there are no inputs, all three quantities are
zero. Membership in the union means that some catalogue input has $G$ as a forward product, which
is exactly the stated intersection condition.
\end{proof}

\begin{definition}[Strong catalogue escape]
\label{def:catalogue-escape}
An exactly assemblable graph strongly escapes $\mathcal C$ under $c$ when every exact input uses
at least one out-of-catalogue component. Its indicator is
\begin{equation}
 \operatorname{Escape}^\star(G;c,\mathcal C)
 =\mathbf1\!\left[\mathcal E_c^\star(G)\neq\varnothing,\quad
 \mathcal E_c^\star(G)\cap\mathfrak C_c(\mathcal C)=\varnothing\right].
\end{equation}
\end{definition}
A catalogue assembler whose outputs lie in $\mathcal G_{\rm cat}(c,\mathcal C)$ therefore has
zero strong-escape yield. For a verified exact-L1 product, strong escape is equivalent to
$G\notin\mathcal G_{\rm cat}(c,\mathcal C)$. This can be checked against a complete canonical
product set or by exhaustive exclusion of all catalogue-only decompositions. A bounded reverse
search cannot establish that exclusion from one novel witness.

The reported experiments did not construct a complete common-reference reachable product set.
Their component-novelty metric concerns the selected verified input. Whole-graph generation removes
catalogue membership from the generative representation; that architectural property alone proves
neither strong escape nor positive probability of any particular out-of-catalogue graph.

\subsection{Conditional flow and ideal population recovery}
\label{app:flow-proof}

We prove Theorem~\ref{thm:ideal-terminal-law} for the reference objective in
Equation~\ref{eq:primary-loss}. Let $X_\star=\operatorname{Enc}_P(G,\alpha)$ denote a clean training
endpoint, and distinguish it from the time-indexed process $X_t$. The training-path measure is
\begin{equation}
 \nu(dG,d\alpha,dP,dt,dx)
 =\mu_{\rm bal}(dG,d\alpha,dP)\,\mathcal U_{(0,1)}(dt)\,
 q_{t\mid1}^P(dx\mid\operatorname{Enc}_P(G,\alpha)).
\end{equation}
For fixed $P$, the encoded endpoint law $q_1^P$ induces
\begin{equation}
 q_t^P(x)=\sum_{x_1}q_1^P(x_1)q_{t\mid1}^P(x\mid x_1),\qquad
 \rho_i(a\mid x,t,P)=\Pr_\nu[X_{\star,i}=a\mid X_t=x,t,P].
 \label{eq:population-path}
\end{equation}
Full source support on the finite allowed alphabets makes $q_t^P(x)>0$ for every $x\in\Omega_P$
and $t<1$. At $t=0$ the path is $q_0^P$; at $t=1$ it is $q_1^P$.

\paragraph{Conditional and marginal rates.}
For $i\in\mathcal V(P)$ and $a,u,v\in\mathcal X_i(P)$, write
$s_{t,i}^P(u\mid a)=(1-t)q_{0,i}^P(u)+t\mathbf1\{u=a\}$ and
$\dot s_i^P(u\mid a)=\mathbf1\{u=a\}-q_{0,i}^P(u)$.
The conditional rate used in the ideal analysis is the $R^*$ construction of
\citet{qin2025defog}, restricted to the allowed alphabet:
\begin{equation}
 R^*_{t,i}(u,v\mid a,P)
 =\frac{[\dot s_i^P(v\mid a)-\dot s_i^P(u\mid a)]_+}
 {|\mathcal X_i(P)|\,s_{t,i}^P(u\mid a)},\qquad u\neq v,\quad 0\leq t<1.
 \label{eq:ideal-conditional-rate}
\end{equation}
Diagonal entries are the negative row sums. Denominators are positive for $t<1$; no value at $t=1$
is needed. Fixed coordinates have zero rates. For $x\neq y$ in $\Omega_P$, the joint conditional
and predicted rates are
\begin{align}
 R_t^*(x,y\mid x_1,P)
 &=\sum_{i\in\mathcal V(P)}\mathbf1\{y_{-i}=x_{-i}\}
 R^*_{t,i}(x_i,y_i\mid x_{1,i},P),\\
 R_t^\theta(x,y\mid P)
 &=\sum_{i\in\mathcal V(P)}\mathbf1\{y_{-i}=x_{-i}\}
 R_{t,i}^\theta(x_i,y_i\mid x,P).
 \label{eq:joint-local-rate}
\end{align}
Here $x_{-i}$ omits coordinate $i$, and Equation~\ref{eq:posterior-rate} defines the predicted
coordinate rate. Diagonal joint rates again make each row sum to zero. States differing at two or
more coordinates have zero instantaneous transition rate.

\begin{proposition}[Program-fixed invariance]
\label{prop:fixed-invariance}
Fix $P\in\mathcal P_{\rm adm}$ and an endpoint in $\Omega_P$. Corruption by
Equation~\ref{eq:conditional-path} stays in $\Omega_P$. The ideal CTMC started from $q_0^P$ stays
there for $t<1$ whenever its coordinate predictions are probability distributions on the allowed
alphabets. Numerical sampling in Equation~\ref{eq:numerical-sampling}, started from $q_0^P$
and using the allowed coordinate alphabets, satisfies
$\Pr[Z_k\in\Omega_P\text{ for all }0\leq k\leq L_{\rm samp}]=1$.
\end{proposition}
\begin{proof}
The corruption law has zero mass outside $\Omega_P$, and the ideal generator has no transition to
its complement. Each numerical coordinate draw stays in its allowed alphabet and $\Pi_P$ restores
the prescribed fixed values. Induction on $k$ proves the discrete statement. These arguments do not
require an accurate denoiser and do not imply molecular validity.
\end{proof}

\begin{proposition}[Posterior minimizer of the reference loss]
\label{prop:masked-bayes}
Under the reference training measure $\nu$, assume unrestricted measurable coordinate predictors
and positive weights $w_i(P)$ for each variable coordinate. Every population minimizer of
Equation~\ref{eq:primary-loss} satisfies
$d_{\theta,i}(a\mid x,t,P)=\rho_i(a\mid x,t,P)$ for all $a\in\mathcal X_i(P)$,
$\nu$-almost surely on contexts with $i\in\mathcal V(P)$.
\end{proposition}
\begin{proof}
Condition on $(X_t,t,P)$ and a variable coordinate. Since $w_i(P)$ is positive and does not depend
on its clean target after this conditioning, its conditional loss is
\begin{equation}
 w_i(P)\left[H(\rho_i)+D_{\rm KL}(\rho_i\Vert d_{\theta,i})\right].
\end{equation}
The entropy is independent of the predictor, and the divergence is nonnegative with equality only
at $d_{\theta,i}=\rho_i$. An unrestricted class can attain all these pointwise minima together.
A population minimizer must attain them almost surely, since a positive-measure excess would
increase the integrated loss. Fixed coordinates contribute no term.
\end{proof}

\begin{corollary}[Population recovery of the local rate]
\label{cor:local-rate-recovery}
For the reference population minimizer in Proposition~\ref{prop:masked-bayes},
$R_t^{\theta^\star}(x,y\mid P)$ equals the conditional expectation of
$R_t^*(x,y\mid X_\star,P)$ given $(X_t=x,t,P)$, almost surely under the path measure.
\end{corollary}
\begin{proof}
Each local conditional rate depends on the endpoint only through its coordinate $X_{\star,i}$.
Substituting the posterior marginal into Equation~\ref{eq:posterior-rate} therefore computes the
required conditional expectation. Summing over coordinates gives the joint rate.
\end{proof}

\begin{proof}[Proof of Theorem~\ref{thm:ideal-terminal-law}]
We verify the forward equation first for one coordinate, then for the conditional product path,
and finally after averaging over endpoints. Fix $P$, $x_1$ and $i$; suppress these indices in
$s_t$ and $\dot s$. For each $v$ and $t<1$, the net flux of Equation~\ref{eq:ideal-conditional-rate}
is
\begin{align}
 \sum_u s_t(u)R_t^*(u,v)
 &=\frac{1}{|\mathcal X_i(P)|}\sum_u
 \bigl([\dot s(v)-\dot s(u)]_+-[\dot s(u)-\dot s(v)]_+\bigr)\\
 &=\dot s(v)-\frac{1}{|\mathcal X_i(P)|}\sum_u\dot s(u)
 =\dot s(v).
\end{align}
The last equality uses $\sum_u\dot s(u)=0$. The product rule then gives the conditional joint
forward equation because the generator changes one coordinate at a time. Averaging over $q_1^P$
and using Bayes' rule yields
\begin{align}
 \partial_t q_t^P(y)
 &=\sum_{x,x_1}q_1^P(x_1)q_{t\mid1}^P(x\mid x_1)R_t^*(x,y\mid x_1,P)\\
 &=\sum_x q_t^P(x)\,
 \mathbb E[R_t^*(x,y\mid X_\star,P)\mid X_t=x,t,P].
\end{align}
Corollary~\ref{cor:local-rate-recovery} identifies the last conditional expectation with the
learned ideal rate for almost every $t$. On each interval $[0,1-\eta]$, $0<\eta<1$, source
positivity bounds all denominators away from zero. The rates are bounded measurable functions on a
finite state space, so the forward equation has a unique solution for initial law $q_0^P$.
Since $q_t^P$ is positive throughout $\Omega_P$, almost-sure equality under the path measure implies
rate equality at every state for Lebesgue-almost every $t$. Changes on the exceptional times do not
change the integral forward equation. The exactly simulated process therefore has law $q_t^P$.

The finite product path converges to the endpoint point mass as $t\uparrow1$. Summing over its
finite endpoint support gives $q_t^P\to q_1^P$ in total variation. This establishes a limit of
marginal laws; a pathwise value at the singular endpoint is not required. Exact decode round trip
maps each endpoint to its source constitution, proving
$\operatorname{Dec}_{\#}q_1^P=\mu_G^P$. The assumed endpoint validity and verified L1 give the
probability-one conclusion. If there are no variable coordinates, the compatible state space is a
singleton and the same conclusions hold for the constant process.
\end{proof}

\begin{corollary}[Decoder contraction and failure probability]
\label{cor:decoder-contraction}
Let $q_1^P$ satisfy the endpoint assumptions of Theorem~\ref{thm:ideal-terminal-law}. Let
$\widetilde q^P$ be any law on
$\overline{\mathfrak X}_{N,K_{\rm cl}}\cup\{\bot_{\rm term}\}$, with $q_1^P$ extended by zero
mass at $\bot_{\rm term}$ and $\operatorname{Dec}(\bot_{\rm term})=\bot_{\rm dec}$.
Define $\mathcal B_P$ as states whose decode fails, is molecularly invalid or fails verified L1.
Then
\begin{equation}
 \|\operatorname{Dec}_{\#}\widetilde q^P-\mu_G^P\|_{\rm TV}
 \leq\|\widetilde q^P-q_1^P\|_{\rm TV},\qquad
 \widetilde q^P(\mathcal B_P)\leq\|\widetilde q^P-q_1^P\|_{\rm TV},
\end{equation}
where $\|p-q\|_{\rm TV}=\frac12\sum_x|p(x)-q(x)|$.
\end{corollary}
\begin{proof}
Every decoder-output event has a preimage in the raw space, so its probability difference is bounded
by the raw-state total variation. Taking the supremum over output events proves the first bound.
The endpoint assumptions give $q_1^P(\mathcal B_P)=0$; applying the same event bound to
$\mathcal B_P$ proves the second.
\end{proof}
The corollary can be applied to the law \emph{after} numerical sampling and terminal readout.
It does not estimate that law's distance from $q_1^P$. In particular, error before the
learned terminal readout cannot be substituted without controlling the readout's effect.

\subsection{Accepted counts under a fixed sampling budget}

Fix a trained model, requested chemistry $c$ and its program-sampling policy. An attempt outcome $Y$
contains its sampled program $P$ and decoded graph or typed failure. Let $Q$ be the common law of
these outcomes and define
\begin{equation}
 A(Y)=\mathbf1\{G\text{ is valid and }\widehat{\mathrm{L1}}(G,c)=1\},\qquad
 z_Q=\mathbb E_Q[A(Y)],
\end{equation}
with $A(Y)=0$ for failures. This definition permits a fresh morphology-specific $P$ on every
attempt. A comparison with a single fixed $P$ is the corresponding special case.

\begin{proposition}[Exact law of matched-budget post-hoc filtering]
\label{prop:posthoc-budget}
For an integer $B\geq0$ and independent attempts with common law $Q$, the accepted-attempt count
$N_B=\sum_{b=1}^BA(Y_b)$ is binomial with parameters $(B,z_Q)$. For all $z_Q\in[0,1]$,
\begin{equation}
 \mathbb E[N_B]=Bz_Q,\qquad
 \operatorname{Var}(N_B)=Bz_Q(1-z_Q),\qquad
 \Pr[N_B=0]=(1-z_Q)^B.
\end{equation}
For $B=0$ the last probability is one, including when $z_Q=1$. If $z_Q>0$, the conditional law of
an accepted attempt is $Q_{\rm post}(y)=Q(y)A(y)/z_Q$, and the number $T_m$ of attempts needed for
$m\geq1$ acceptances satisfies $\mathbb E[T_m]=m/z_Q$.
For two proposal laws under the same attempt budget,
$\mathbb E[N_B^{(1)}-N_B^{(0)}]=B(z_{Q_1}-z_{Q_0})$.
\end{proposition}
\begin{proof}
Each acceptance indicator is Bernoulli with parameter $z_Q$, and independence makes their sum
binomial. Its expectation and variance give the first two identities; all $B$ attempts fail with
probability $(1-z_Q)^B$. For $z_Q>0$, conditioning on acceptance gives the normalized law by
Bayes' rule. The waiting times between acceptances are independent geometric variables on
$\{1,2,\ldots\}$ with mean $1/z_Q$; summing $m$ of them proves the waiting-time formula.
Subtracting the expected counts proves the comparison identity, without requiring independence
between the two arms.
\end{proof}
If $z_Q=0$, accepted outcomes have no conditional law and $T_m=\infty$ almost surely for $m\geq1$.
The proposition counts accepted attempts, including duplicates. It does not establish that program
conditioning increases acceptance; that is an empirical comparison. If requested programs are fixed
in unequal strata, independent acceptance indicators generally have different probabilities and
their sum is Poisson-binomial. Retraining variation is separate from both sampling laws.

\subsection{Route termination and evidence completeness}

\begin{lemma}[Termination of bounded route assessment]
\label{lem:planner-termination}
Fix a route contract $\mathfrak B$ with finite maximum depth $D_{\mathfrak B}$. Assume finitely many
proposed expansions per molecule and reactants per expansion, a terminating evidence and
forward-verification call for each proposal, and a search that assesses each proposed expansion
finitely many times. For an admitted root trace, recursive assessment terminates with a dossier
or typed failure in $\mathfrak D\cup\mathfrak F_{\mathfrak B}$.
\end{lemma}
\begin{proof}
Induct on remaining depth. At depth zero, the terminating availability query either admits a
terminal or returns a failure without further expansion. At positive remaining depth, only finitely
many proposals are assessed, each using terminating local calls. Every proposed child has one
fewer remaining level, so its assessment terminates by induction. A finite number of finite child
assessments and proposals therefore terminates at the parent. The same argument applies if a
finite outer root loop is used. Rejecting ancestor identities prevents circular dependencies from
being admitted; the finite depth and no-infinite-retry assumptions establish termination.
\end{proof}

\begin{invariant}[Synthesis-dossier admission]
\label{inv:dossier-soundness}
Under the frozen contract $\mathfrak B$, a successful assessor return $T\in\mathfrak D$ is admitted
only after checking $\mathrm{Complete}_{\mathfrak B}(G,T,P)=1$. Every other outcome is recorded as
a typed failure. This is the admission rule in Algorithm~\ref{alg:forge-inference}, not a theorem
that bounded search finds every supported chemical route.
\end{invariant}
The predicate checks molecular validity, a unique verified root trace, a well-formed finite tree,
evidence admission and forward replay for every preparation step, and availability for every leaf.
Its factors are defined in Appendix~\ref{app:route-assessment}. A failure establishes only that the
assessor did not return a complete dossier under the recorded contract. It does not exclude a route
under a larger budget, other admitted chemistry or another supplier snapshot.

\clearpage

\section{Assembly families and reaction schemes}
\label{app:chemistry-examples}
\label{app:preliminaries}

The assembly schemes depict the chemistry introduced in Section~\ref{sec:reaction-background}.
The 22-family structure atlas pairs products with precursor origins and exact assembly traces
(Figure~\ref{fig:22-atlas}).

\subsection{Assembly schemes}
\label{app:chemistry}
The precursor roles and assembly patterns are summarized in Table~\ref{tab:reaction-family-background}.

\begin{longtable}{@{}>{\raggedright\arraybackslash}p{.36\textwidth}>{\raggedright\arraybackslash}p{.58\textwidth}@{}}
\caption{\textbf{Registered assembly families and precursor roles.} Source role names are retained;
a family can contain multiple admitted program variants and role aliases. A list of role names
does not specify that every role occurs in a single reaction. The atom-mapped program determines
the participating components and transformation order.}\label{tab:reaction-family-background}\\
\toprule Family & Registered source roles\\\midrule\endfirsthead
\toprule Family & Registered source roles (continued)\\\midrule\endhead
\bottomrule\endfoot
A3 amine--aldehyde--alkyne & aldehyde, alkyne, amine head \\
Acid--epoxide diester & amine acid head, epoxide, hydrophobic acid \\
AEMA & AEMA, amine core, thiol periphery \\
Aldehyde Ugi 3-CR & aldehyde side, amine side, isocyanide side \\
Aldehyde Ugi 4-CR & aldehyde, amine, carboxylic acid, isocyanide \\
$\alpha$-Isocyanoester dihydroimidazole & amine head, coupled ketone, isocyanide \\
Amine alkylation & amine head, bromoester arm \\
Amine--epoxide & amine head, epoxide tail \\
Aryl reductive amination & amine head, coupled aldehyde \\
Aza-Michael acrylamide & acceptor tail, amine head \\
Aza-Michael acrylate & acceptor tail, acrylate, alkyl acrylate, amine, amine head, amine nucleophile, amine precursor, conjugated carbonyl acceptor \\
Disulfide Michael & amine head, disulfide acceptor \\
Epoxide O-acylation & acyl tail, amine head, epoxide tail \\
iPhos & amine head, phosphate tail \\
Ketone--isocyanide amide & amine head, coupled ketone, isocyanide \\
Ketone Ugi 4-CR & amine, amine head, carboxylic acid, coupled ketone, isocyanide, ketone \\
Maleate & amine head, maleate \\
O-esterification & acid tail, aminoalcohol head \\
Passerini & aldehyde tail, amine acid head, isocyanide tail \\
Preassembled thiol-yne & alkynoic linker, amine head, thiol tail \\
Reductive amination & amine head, coupled aldehyde \\
Thiolactone & acrylate tail, amine head, thiolactone region \\
\end{longtable}

\subsection{Reaction schemes across the assembly families}
\label{app:all-family-schemes}
Figures~\ref{fig:reaction-schemes}--\ref{fig:rxn-ordered-d} show the registered transformations
for all 22 families. The companion schemes use exact-replay-verified representative assemblies.
Unchanged hydrocarbon fragments are abbreviated by $R$ labels local to each arrow; heteroatoms,
ring systems and the connecting molecular framework are retained. Ordered stages use the preceding
intermediate after the indicated repeated incorporations. Auxiliary reagents, salts, work-up and
stereochemistry absent from the constitutional record are omitted. The drawings specify
reaction-enumerated support without assigning experimental outcomes to the depicted combinations.
Precursor-origin boundaries can differ from the functional headgroup, linker and tail regions
introduced in Section~\ref{sec:lipid-background}. An L1 pass requires forward replay on the
recovered precursors to reproduce the complete constitutional lipid graph exactly.

% [tbp] rather than [H]: the schemes are too tall to sit inline, and forcing them here broke the
% page early and left LaTeX stretching the section glue above to fill it.
\begin{figure}[tbp]
\centering
\input{figures/forge_reaction_schemes/schemes_body.tex}
\caption{\textbf{Representative assembly chemistries.} Monochrome skeletal drawings show the
reactants and products of each registered assembly. Hydrogens are drawn only on the atoms whose substitution decides
whether a product can be decomposed back into precursors: the head N--H, and the
$\alpha$-carbon C--H that the Ugi reaction takes from the aldehyde. $R^i$ stand for the
source-linked components of the frozen library, which are larger than the fragments drawn here.
The Ugi variant has no carboxylic-acid reactant component.}
\label{fig:reaction-schemes}
\end{figure}

\clearpage


\clearpage
\begin{figure}[H]\centering
\includegraphics[width=\linewidth,height=.80\textheight,keepaspectratio]{figures/all_family_reactions/multicomponent_a.pdf}
\caption{\textbf{Multicomponent assembly: A3 and Ugi four-component reactions.}
\textbf{a}, A3 coupling creates C--N and C--C bonds while retaining the alkyne.
\textbf{b}, Aldehyde Ugi 4-CR combines an amine, aldehyde, isocyanide and carboxylic acid.
The acid supplies the N-acyl unit in the four-component reaction.}
\label{fig:rxn-multicomponent-a}\end{figure}

\clearpage
\begin{figure}[H]\centering
\includegraphics[width=\linewidth,height=.80\textheight,keepaspectratio]{figures/all_family_reactions/ketone_ugi4.pdf}
\caption{\textbf{Ketone Ugi four-component assembly.} The ketone reacts with an amine, isocyanide
and carboxylic acid. Both carbon substituents remain attached to the ketone-derived center; the
preassembled ester-bearing arms are retained, and the acid supplies the N-acyl unit.}
\label{fig:rxn-ketone-ugi4}\end{figure}

\clearpage
\begin{figure}[H]\centering
\includegraphics[width=\linewidth,height=.80\textheight,keepaspectratio]{figures/all_family_reactions/multicomponent_b.pdf}
\caption{\textbf{Multicomponent assembly: aminoamide, dihydroimidazole and Passerini reactions.}
\textbf{a}, The ketone--isocyanide branch forms an aminoamide from three components.
\textbf{b}, An $\alpha$-isocyanoester participates in dihydroimidazole ring formation, retaining its ester.
\textbf{c}, Passerini assembly produces an $\alpha$-acyloxy amide from an aldehyde, acid and isocyanide;
the acid-derived acyl group attaches through oxygen.}
\label{fig:rxn-multicomponent-b}\end{figure}

\clearpage
\begin{figure}[H]\centering
\includegraphics[width=\linewidth,height=.80\textheight,keepaspectratio]{figures/all_family_reactions/conjugate_additions.pdf}
\caption{\textbf{Conjugate-addition families.}
\textbf{a}, Amine addition to an acrylamide retains the amide interface.
\textbf{b}, Addition to a disulfide-containing acrylate retains its S--S bond.
\textbf{c}, The illustrated maleate program uses the N-addition branch and retains both ester groups.
For repeated programs, one event is shown and the incorporated-event count is indicated.
The aza-Michael acrylate reaction is shown in Figure~\ref{fig:reaction-schemes}.}
\label{fig:rxn-conjugate}\end{figure}

\clearpage
\begin{figure}[H]\centering
\includegraphics[width=\linewidth,height=.80\textheight,keepaspectratio]{figures/all_family_reactions/functionalization.pdf}
\caption{\textbf{Nitrogen functionalization.}
\textbf{a}, Aryl reductive amination combines condensation with reduction.
\textbf{b}, N-alkylation replaces the alkyl bromide leaving group.
\textbf{c}, Amine attack opens an epoxide and retains the resulting alcohol.
Repeated incorporations are labeled.}
\label{fig:rxn-functionalization}\end{figure}

\clearpage
\begin{figure}[H]\centering
\includegraphics[width=\linewidth,height=.80\textheight,keepaspectratio]{figures/all_family_reactions/oxygen_phosphate.pdf}
\caption{\textbf{Esterification and phosphate opening.}
\textbf{a}, O-esterification joins an aminoalcohol to an acid while retaining the basic nitrogen.
\textbf{b}, iPhos assembly opens a C--O bond of a cyclic phosphate to form the N--CH$_2$CH$_2$--O--P
connection. Its deprotonated product representation is preserved.}
\label{fig:rxn-oxygen-phosphate}\end{figure}

\clearpage
\begin{figure}[H]\centering
\includegraphics[width=\linewidth,height=.80\textheight,keepaspectratio]{figures/all_family_reactions/ordered_assemblies_a.pdf}
\caption{\textbf{Ordered assembly through hydroxy-ester and dual-acceptor intermediates.}
\textbf{a}, Acid--epoxide opening first produces a hydroxy ester; subsequent esterification attaches
the aminoacid head through its carboxyl group.
\textbf{b}, The AEMA program consumes the acrylate in an aza-Michael addition, then the remaining
methacrylate in a thiol addition. Each second-stage reactant is the corresponding first-stage
intermediate; $R$ abbreviations are local to each arrow.}
\label{fig:rxn-ordered-a}\end{figure}

\clearpage
\begin{figure}[H]\centering
\includegraphics[width=\linewidth,height=.80\textheight,keepaspectratio]{figures/all_family_reactions/ordered_assemblies_b.pdf}
\caption{\textbf{Epoxide opening followed by O-acylation.}
Amine--epoxide additions first form aminoalcohol arms. The resulting alcohols then react with acyl
chlorides to form ester linkages. The representative program incorporates two epoxides and two
acyl groups; one event from each stage is drawn, with its multiplicity indicated.}
\label{fig:rxn-ordered-b}\end{figure}

\clearpage
\begin{figure}[H]\centering
\includegraphics[width=\linewidth,height=.80\textheight,keepaspectratio]{figures/all_family_reactions/ordered_assemblies_c.pdf}
\caption{\textbf{Thiol--yne tail assembly followed by amidation.}
Two identical thiols ($R_1=R_2$ within each arrow) add across a terminal alkyne to form the
bis(thioether) acid, which then undergoes amidation. The first arrow denotes the net double addition.}
\label{fig:rxn-ordered-c}\end{figure}

\clearpage
\begin{figure}[H]\centering
\includegraphics[width=\linewidth,height=.80\textheight,keepaspectratio]{figures/all_family_reactions/ordered_assemblies_d.pdf}
\caption{\textbf{Thiolactone aminolysis followed by thiol--Michael addition.}
Aminolysis opens the thiolactone to form an amide and release a thiol. The subsequent addition
to an acrylate forms a thioether. The retained sulfur framework distinguishes this transformation
from disulfide formation.}
\label{fig:rxn-ordered-d}\end{figure}


\end{document}
